% concatenated from Almost_Golomb_V3.tex
\documentclass[12pt]{article}
\usepackage[utf8]{inputenc}
\usepackage{amsmath,amssymb,amsthm}
\usepackage[margin=2.5cm]{geometry}
\usepackage[dvipsnames]{xcolor}
\usepackage{hyperref}
\hypersetup{
  colorlinks=true,
  linkcolor=MidnightBlue,
  citecolor=MidnightBlue,
  urlcolor=MidnightBlue,
  pdftitle={Almost Golomb Sequences},
  pdfauthor={Beno\^it Cloitre}
}
\usepackage{enumitem}

\newtheorem{theorem}{Theorem}[section]
\newtheorem{lemma}[theorem]{Lemma}
\newtheorem{proposition}[theorem]{Proposition}
\newtheorem{corollary}[theorem]{Corollary}
\theoremstyle{definition}
\newtheorem{definition}[theorem]{Definition}
\newtheorem{example}[theorem]{Example}
\theoremstyle{remark}
\newtheorem{conjecture}[theorem]{Conjecture}
\newtheorem{remark}[theorem]{Remark}

\newcommand{\N}{\mathbb{N}}
\newcommand{\Z}{\mathbb{Z}}
\newcommand{\Ical}{\mathcal{I}}
\newcommand{\one}{\mathbf{1}}
\newcommand{\oeis}[1]{\href{https://oeis.org/#1}{\textbf{#1}}}

\title{Almost Golomb Sequences}
\author{Beno\^it Cloitre}
\date{}

\hfuzz=2pt
\begin{document}
\maketitle

\begin{abstract}
Golomb's sequence is the unique nondecreasing sequence of positive integers
in which each $n$ appears exactly $a(n)$ times. It satisfies the global
self-referential rule
\[
  a\bigl(a(n)+a(n-1)+\cdots+a(1)\bigr)=n,
\]
grows smoothly like a power of $n$ governed by the golden ratio,
and is not $k$-regular for any $k\ge 2$.

We introduce almost Golomb sequences, obtained by truncating the
cumulative sum to a sliding window of fixed size $r$,
\[
  a\bigl(a(n)+a(n-1)+\cdots+a(n-r+1)\bigr)=n.
\]
This finite-memory truncation changes the nature of the sequence completely.
The smooth power law gives way to oscillatory linear growth, and the
sequence becomes $r$-regular for every $r\ge 2$. For small values of $r$
we establish explicit denesting formulas, prove that $a(n)/n$ does not
converge, and uncover combinatorial structure including a cellular
automaton and a palindromic substitution.

A numerical surprise emerges when one varies $r$. The maximum multiplicity
across the family of sequences is governed by Golomb's sequence itself.
The sequence that was truncated reappears as the law controlling the
family it generated.
\end{abstract}

\noindent\textbf{2020 Mathematics Subject Classification:}
11B85 (primary), 11B37, 68R15 (secondary).

\noindent\textbf{Keywords:} self-referential sequences, Golomb's sequence,
$k$-regular sequences, automatic sequences, sliding-window recurrences,
denesting, non-convergence.

%-------------------------------------------------------------------
\section{Introduction}
%-------------------------------------------------------------------

\subsection{Golomb's sequence}\label{sec:golomb}

Golomb's sequence $(a(n))_{n\ge 1}$ (OEIS \oeis{A001462},~\cite{OEIS}) is the unique nondecreasing
sequence of positive integers in which the integer $n$ appears exactly $a(n)$ times.
It begins
\[
  1, 2, 2, 3, 3, 4, 4, 4, 5, 5, 5, 6, 6, 6, 6, 7, 7, 7, 7, \ldots
\]
In particular, $a$ satisfies
\begin{equation}\label{eq:golomb}
  a\!\left(\sum_{k=1}^{n} a(k)\right) = n, \qquad n\ge 1,
\end{equation}
a formulation recorded in the OEIS entry \oeis{A001462}.
Golomb's sequence is one of the most studied self-referential sequences in the
world of integer sequences, introduced by Golomb~\cite{Golomb66} and prominently
featured by Sloane among his favourite sequences.
Its asymptotic behavior is smooth:
\begin{equation}\label{eq:golomb-asymp}
  a(n) \sim \varphi^{2-\varphi}\, n^{\varphi-1}, \qquad n\to\infty,
\end{equation}
where $\varphi=(1+\sqrt{5})/2$ is the golden ratio~\cite{petermann95},
and $f\sim g$ means $f(n)/g(n)\to 1$.
Since $a(n)$ grows like $n^{\varphi-1}$ with irrational algebraic exponent
$\varphi-1\approx 0.618$, Golomb's sequence is not $k$-regular for any integer
$k\ge 2$: if it were, its subsequences along geometric progressions
$(a(k^m))_{m\ge 0}$ would satisfy a linear recurrence with integer
coefficients~\cite[Theorem~2.2]{AS92}, requiring the asymptotic growth factor
$k^{\varphi-1}$ to be an algebraic integer; but $k^{\varphi-1}$ is
transcendental by the Gelfond--Schneider theorem (which states that $\alpha^\beta$
is transcendental whenever $\alpha$ is algebraic, $\alpha\notin\{0,1\}$, and $\beta$
is an irrational algebraic number), a contradiction.

The combinatorial depth of Golomb's sequence has recently been made explicit by
Claes and Miyamoto~\cite{CM26}, who define the \emph{golombic operator}
$\Gamma$: given a sequence $a$, $\Gamma(a)$ is the sequence whose $n$-th
run has length $a(n)$ (that is, value $n$ is repeated $a(n)$ times).
They show that Golomb's sequence is the unique fixed point of $\Gamma$
in the class of nondecreasing sequences, establishing a rich algebraic
structure. Our almost Golomb sequences occupy a different regime. They are not fixed
points of $\Gamma$, but their run lengths are bounded (depending on $r$) and the
corresponding run-length sequence is $r$-automatic. For $r=2$ in particular the
run lengths lie in $\{1,2\}$ (Proposition~\ref{prop:N2_multiplicity}).
Numerically, the iterates $\Gamma^k(a_r)$ appear to converge to $G$ in the
prefix topology, suggesting a deeper dynamical link that we leave for future work.

\subsection{Almost Golomb sequences}

What happens when one replaces the full cumulative sum in
\eqref{eq:golomb} by a sliding window of the $r$ most recent terms. This leads to
sequences that we call \emph{almost Golomb sequences}: they share the meta-recursive
spirit of Golomb's sequence~\cite{Hofstadter}, but are built from a finite-memory rule.

\begin{definition}\label{def:almost-golomb}
  Let $r\ge 2$. The \emph{almost Golomb sequence of order $r$} is the
  nondecreasing sequence of positive integers $(a(n))_{n\ge 1}$,
  extended by $a(k)=0$ for $k\le 0$, satisfying
  \begin{equation}\label{eq:def}
    a\bigl(a(n)+a(n-1)+\cdots+a(n-r+1)\bigr) = n, \qquad n\ge 1,
  \end{equation}
  with $a(n)$ minimal: for every $n\ge 2$, no integer $m$ with
  $a(n-1)\le m<a(n)$ satisfies
  $a(m + a(n-1)+\cdots+a(n-r+1)) = n$.
\end{definition}

\begin{remark}
This is an implicit definition: the constraint~\eqref{eq:def}
involves the value of $a$ at the position $a(n)+\cdots+a(n-r+1)$,
which in general lies beyond~$n$.
The existence and uniqueness of the sequence, as well as the
consistency of the greedy rule, are established in
Section~\ref{sec:general} (Theorem~\ref{thm:global_run}).
\end{remark}

The almost Golomb sequences belong to a broader family of monotone self-referential
sequences admitting an automatic structure via denesting. A comparative table
of related sequences is collected in Appendix~\ref{app:nested}.

This truncation from an infinite to a finite memory window changes the nature of the sequence.
Golomb's global rule forces a smooth, sublinear growth governed by an irrational
exponent. The finite-window rule instead imposes linear growth accompanied by
discrete, $r$-automatic oscillations.
The sequence transitions from the realm of analytic power laws into the realm of
automatic sequences.

A sequence is \emph{$r$-regular} (in the sense of Allouche and Shallit~\cite{AS92})
if all its subsequences of the form $(f(r^i n+j))_{n\ge 0}$ lie in a finitely
generated $\mathbb{Z}$-module. Informally, its values along arithmetic progressions
in base $r$ satisfy linear recurrences.
This framework is well adapted to sequences defined by divide-and-conquer rules.

The paper is organised as follows.
Section~\ref{sec:general} establishes $r$-regularity universally.
Sections~\ref{sec:r2}--\ref{sec:r3} give complete constructive proofs
for $r\in\{2,3\}$.
Sections~\ref{sec:r4}--\ref{sec:r5} treat $r\in\{4,5\}$ via explicit
recurrences and Walnut certification.
Section~\ref{sec:nonconv} proves non-convergence of $a(n)/n$.
Section~\ref{sec:combinatorial} presents combinatorial interpretations
for $r\in\{2,3\}$.
Sections~\ref{sec:golomb_return}--\ref{sec:open} present the Golomb
meta-structure conjectures and open questions.

The main contributions of this paper are:
\begin{enumerate}
  \item A proof (Theorem~\ref{thm:structural_regularity}, Section~\ref{sec:general}) that $(a(n))$ is $r$-regular
        for every $r\ge 2$, by reducing the automaticity of the first-difference
        sequence $d(n)=a(n+1)-a(n)$ to a finite window of local data.
  \item Explicit denesting of the almost Golomb sequence for each
        $r\in\{2,3,4,5\}$ into $r$-adic divide-and-conquer formulas
        (Sections~\ref{sec:r2}--\ref{sec:r5}, Theorems~\ref{thm:r2}, \ref{thm:r3}, \ref{thm:r4}, \ref{thm:r5}).
        The cases $r\in\{2,3\}$ are proved directly; the cases $r\in\{4,5\}$
        use a hybrid method combining explicit recurrence identification
        with automatic certification by Walnut~\cite{Shallit2022}.
        The non-generic terms are controlled by explicit $r$-automatic
        binary correction sequences: one for $r=3$, four
        for $r=4$, and a disjoint pair $(\varepsilon,\eta)$ for $r=5$.
  \item For $r=2$, a Mallows-type nested recurrence
        (Proposition~\ref{prop:mallows-r2}):
        \[
          a(n+1)=1+a\bigl(n+1-a(a(n)+1)+a(a(n)-1)\bigr),
        \]
        analogous to the Mallows recurrence $G(n+1)=1+G(n+1-G(G(n)))$
        for Golomb's sequence.
  \item For each $r\in\{2,3,4,5\}$, a proof that $a(n)/n$ does not converge,
        identifying two explicit rational limit points
        (Section~\ref{sec:nonconv}, Theorems~\ref{thm:ratio2}, \ref{thm:ratio3},
        \ref{thm:ratio4}, \ref{thm:ratio5}).
        For $r\in\{4,5\}$ the proofs use
        Lemma~\ref{lem:auto-geometric} to extract the periodic
        corrector pattern along geometric progressions from the DFAO.
  \item Combinatorial interpretations for the cases $r=2$ and $r=3$
        (Section~\ref{sec:combinatorial}): a second-bit characterization of
        multiplicities and a Boolean cellular automaton for $r=2$,
        and a palindromic substitution rule for $r=3$.
  \item A conjectural Golomb meta-structure (Section~\ref{sec:golomb_return}):
        numerical evidence suggests that the maximal multiplicity
        $M(r)=\sup_n N_r(n)$ is governed by Golomb's sequence.
        More precisely, if $j_k$ denotes the smallest order $r$ such that
        $M(r)\ge k$, then the data support the identities
        \[
          j_{k+1}-j_k = G(k), \qquad j_k = S(k-1)+2,
        \]
        where $G$ is Golomb's sequence and $S(k)=\sum_{n\le k}G(n)$.
        We reduce this law to two precise boundary conjectures
        (the Prefix Conjecture and the Domination Lemma).
\end{enumerate}

%-------------------------------------------------------------------
%-------------------------------------------------------------------
\section{Universal $r$-regularity}\label{sec:general}
%-------------------------------------------------------------------

In this section we prove that the almost Golomb sequence of order $r$ is
$r$-regular for every $r\ge 2$. The proof has two parts.

First, we establish the global run structure of the sequence: the values are
not skipped, the increments are always $0$ or $1$, and for every $m\ge 3$
the first occurrence of $m$ is at position
\[
  S_m^{(r)}:=a(m)+a(m-1)+\cdots+a(m-r+1).
\]
Second, we use this run structure to express the
difference sequence at scale $rn+i$ in terms of a finite local window,
and then apply the Allouche--Shallit criterion (Theorem~\ref{thm:AS}).

\subsection{Global run structure}

We extend the sequence by setting $a(k)=0$ for $k\le 0$.
For $n\ge 1$, write $S_n:=S_n^{(r)}:=a(n)+\cdots+a(n-r+1)$,
so the defining equation reads $a(S_n)=n$.

\begin{lemma}\label{lem:anchor-growth}
For every $r\ge 2$, the sequence of starting positions $(S_n)_{n\ge 1}$ is strictly increasing.
\end{lemma}

\begin{proof}
If $S_n\le S_{n-1}$ for some $n\ge 2$, then $n=a(S_n)\le a(S_{n-1})=n-1$,
a contradiction.
\end{proof}

The following theorem describes the run structure of $a$.

\begin{theorem}\label{thm:global_run}
For every $r\ge 2$, there exists a unique almost Golomb sequence of order $r$.
It satisfies:
\begin{enumerate}[label=(\roman*)]
  \item $d(n):=a(n+1)-a(n)\in\{0,1\}$ for every $n\ge 1$;
  \item every positive integer occurs at least once;
  \item for every $m\ge 3$, the first occurrence of $m$ is exactly at $S_m$;
  \item for every $m\ge 3$, the run of value $m$ is exactly $[S_m,\;S_{m+1}-1]$.
\end{enumerate}
\end{theorem}

\begin{proof}
We construct a sequence $b$ by its runs, show that it satisfies
Definition~\ref{def:almost-golomb}, and then prove uniqueness.

\textit{Initial prefix.}
For every $r\ge 2$, the constraint~\eqref{eq:def} at $n=1$ gives $a(S_1)=1$
with $S_1=a(1)$, forcing $a(1)=1$.
At $n=2$, minimality gives $a(2)=\cdots=a(S_3-1)=2$
(since $a(k)=0$ for $k\le 0$ makes $S_2=a(2)+a(1)=a(2)+1$,
and the constraint $a(S_2)=2$ is first satisfiable with $a(2)=2$).
This prefix is nondecreasing with increments in $\{0,1\}$, its image is
$\{1,2\}$, and $S_3\ge 4$.
Define $b$ identically on this prefix.

\textit{Inductive run construction.}
The induction hypothesis $\mathcal H_m$ asserts: the prefix of $b$
up to $S_m^b-1$ is nondecreasing with increments in $\{0,1\}$,
its image is exactly $\{1,\ldots,m-1\}$, runs are as stated,
every run has length $\le r$, and $S_m^b\ge m+1$.

Assume $\mathcal H_m$. Set $b(S_m^b):=m$ and define $b\equiv m$ on
$[S_m^b,\,S_{m+1}^b-1]$. We must show $S_{m+1}^b>S_m^b$,
i.e.\ $L_m^b:=b(m+1)-b(m+1-r)\ge 1$.
Since $\mathcal H_m$ gives $S_m^b\ge m+1$, position $m+1$ lies within
a previously defined run, so the value $b(m+1)$ is already determined
at this stage.

Suppose $L_m^b=0$. Then $b$ is constant on $[m+1-r,m+1]$ with value $v$.
If $v=m$: the value $m$ already appears at position $m+1-r$.
Since the first occurrence of $m$ is at $S_m^b$, this implies $m+1-r\ge S_m^b$.
But $\mathcal H_m$ gives $S_m^b\ge m+1$, hence $m+1-r\ge m+1$, impossible since $r\ge 2$.
If $v<m$: the run of $v$ has length $\ge r+1$, contradicting $\mathcal H_m$.
Hence $L_m^b\ge 1$, $S_{m+1}^b>S_m^b$.
Moreover $L_m^b=\sum_{j=0}^{r-1}d(m-j)\le r$ since $d\in\{0,1\}$,
so the run of $m$ has length $\le r$ and $\mathcal H_{m+1}$ holds.

\textit{Verification of the greedy rule.}
For positions $q<S_3^b$ this is checked directly.
Fix a position $q\ge S_3^b$ and let $m:=b(q)$ be the value assigned by
the run construction. We must show that $m$ is the smallest integer
$v\ge b(q-1)$ satisfying $b(v+b(q-1)+\cdots+b(q-r+1))=q$.

For the candidate $v=m$, the window sum is
$T_q(m):=m+b(q-1)+\cdots+b(q-r+1)=S_q^b$, and $b(S_q^b)=q$ since
runs begin at their starting position. Hence $v=m$ satisfies the defining equation.

It remains to show minimality.
If $q=S_m^b$ (start of the run of $m$): then $b(q-1)=m-1$, so $v\ge m-1$.
The candidate $v=m-1$ gives window sum $T_q(m-1)=S_q^b-1<S_q^b$.
Since $S_q^b$ is the starting position for value $q$, we have $b(S_q^b-1)<q$,
so $v=m-1$ fails.
If $q>S_m^b$ (interior of the run): then $b(q-1)=m$, so $v\ge m$ and
$m$ is already the smallest candidate.
Hence $b$ satisfies Definition~\ref{def:almost-golomb}, so existence holds
and $b$ has properties (i)--(iv).

For uniqueness, let $c$ be any sequence satisfying
Definition~\ref{def:almost-golomb}.
The constraint at $n=1,2$ forces $c(1)=b(1)=1$ and $c(n)=b(n)=2$
for $2\le n\le S_3^b-1$, so $c$ and $b$ share the same initial prefix.
Assume inductively that $c$ agrees with $b$ up to $S_m^b-1$.
Then the next admissible value and its maximal run compatible with
Definition~\ref{def:almost-golomb} are forced, so the run of value $m$
in $c$ coincides with that of $b$.
Hence $c=b$ term by term.
\end{proof}

\begin{lemma}\label{lem:runlength-d}
For every $r\ge 2$ and $m\ge 3$,
\[
  L_m:=S_{m+1}-S_m=a(m+1)-a(m+1-r)=\sum_{j=0}^{r-1} d(m-j).
\]
\end{lemma}

\begin{proof}
By Theorem~\ref{thm:global_run}$(iv)$, $L_m=S_{m+1}-S_m$.
Telescoping: $S_{m+1}-S_m=a(m+1)-a(m+1-r)=\sum_{j=0}^{r-1}d(m-j)$.
\end{proof}

Composing the map $n\mapsto S_n$ with itself gives the following identity.

\begin{lemma}\label{lem:nested-anchor}
For every $r\ge 2$ and $n\ge 1$,
\[
  S_{S_n}=rn-R_n, \qquad R_n:=\sum_{j=1}^{r-1}(r-j)\,d(S_n-j)\in[0,\,r(r-1)/2].
\]
\end{lemma}

\begin{proof}
Since $a(S_n)=n$, one has $a(S_n-k)=n-\sum_{j=1}^{k}d(S_n-j)$ for $0\le k\le r-1$.
Summing: $S_{S_n}=rn-\sum_{j=1}^{r-1}(r-j)d(S_n-j)$.
\end{proof}

The value $d(rn+i)$ can be read from a bounded window of local data.

\begin{lemma}\label{lem:local_projection}
Let $A\ge 3$.
\begin{enumerate}[label=(\roman*)]
  \item For $X\ge 0$: $d(S_A+X)=1\iff \sum_{j=0}^{\delta-1}L_{A+j}=X+1$
    for some $\delta\in\{1,\ldots,X+1\}$.
    Hence $d(S_A+X)$ is determined by $(L_A,\ldots,L_{A+X})$.
  \item For $Y\ge 1$: $d(S_A-1-Y)=1\iff \sum_{j=1}^{\delta}L_{A-j}=Y$
    for some $\delta\in\{1,\ldots,Y\}$.
    Hence $d(S_A-1-Y)$ is determined by $(L_{A-1},\ldots,L_{A-Y})$.
    Also, $d(S_A-1)=1$ always.
\end{enumerate}
\end{lemma}

\begin{proof}
By Theorem~\ref{thm:global_run}, $d(k)=1\iff k=S_m-1$ for some $m$.
For $(i)$: $d(S_A+X)=1\iff S_A+X+1=S_{A+\delta}$ for some $\delta\ge 1$,
i.e.\ $\sum_{j=0}^{\delta-1}L_{A+j}=X+1$; since $L_u\ge 1$, one has $\delta\le X+1$.
For $(ii)$: similarly $d(S_A-1-Y)=1\iff S_A-Y=S_{A-\delta}$,
i.e.\ $\sum_{j=1}^{\delta}L_{A-j}=Y$; and $d(S_A-1)=1$ since $S_A-1$ is the last
position of the run of $A$.
\end{proof}

The run structure gives a direct characterization of the value $a(K)$.

\begin{lemma}\label{lem:anchor-characterization}
For every $r\ge 2$, $m\ge 3$, and every integer $K\ge 1$,
\[
  a(K)=m \iff S_m \le K < S_{m+1}.
\]
\end{lemma}

\begin{proof}
By Lemma~\ref{lem:anchor-growth} the positions $S_m$ are strictly increasing,
and by Theorem~\ref{thm:global_run}(iv) the run of value~$m$ is exactly
$[S_m,\,S_{m+1}-1]$.
\end{proof}


\subsection{The Allouche--Shallit criterion}

A sequence $(f(n))_{n\ge 0}$ with values in a finite set is \emph{$r$-automatic}~\cite{Cobham}
if its $r$-kernel $\{(f(r^in+j))_{n\ge 0}:i\ge 0,\,0\le j<r^i\}$ is finite.
A sequence with integer values is \emph{$r$-regular}~\cite{AS92} if the $\mathbb{Z}$-module
spanned by its $r$-kernel is finitely generated. Every $r$-automatic sequence is $r$-regular.

The following criterion underlies all automaticity results in this paper.

\begin{theorem}[\cite{AS12}]\label{thm:AS}
Let $(U(n))_{n\ge 0}$ take values in a finite set, and let $r\ge 2$.
Then $U$ is $r$-automatic if there exist $a,b\ge 0$ and $n_0\ge 0$ such that
for all $n\ge n_0$ and $0\le i<r$,
\[
  U(rn+i) = f_i\bigl(U(n-a),\ldots,U(n+b)\bigr)
\]
for some functions $f_i$.
\end{theorem}

The following observation will be used in Section~\ref{sec:nonconv}
to evaluate automatic sequences along geometric progressions.

\begin{lemma}\label{lem:auto-geometric}
Let $u$ be a $b$-automatic sequence generated by a DFAO reading
base-$b$ expansions from most significant to least significant digit,
and let $P$, $Q$ be fixed base-$b$ words.
Then the sequence $k\mapsto u([P\,0^k\,Q]_b)$ is ultimately periodic.
\end{lemma}

\begin{proof}
Let $(\mathcal Q,\Sigma,\delta,q_0,\tau)$ be the DFAO.
After reading the prefix $P$, the automaton is in a fixed state $q_P$.
Reading $0^k$ applies the map $q\mapsto\delta(q,0)$ exactly $k$ times,
producing a sequence of states that is ultimately periodic
(since the state set is finite).
Reading the suffix $Q$ and applying $\tau$ preserves this periodicity.
\end{proof}

\subsection{Universal $r$-regularity}

Set $d(0):=1$, so that $a(n)=\sum_{k=0}^{n-1}d(k)$ for all $n\ge 1$.

\begin{theorem}\label{thm:structural_regularity}
For every $r\ge 2$, $d(n)$ is $r$-automatic and $(a(n))$ is $r$-regular.
\end{theorem}

\begin{proof}
Set $C:=2r-3$, $D:=r(r+1)/2-1$, $n_0:=\max(r+1,2r-3)$.
Fix $n\ge n_0$ and $0\le i<r$.

\textit{From scale $rn+i$ to scale $S_n$.}
By Lemma~\ref{lem:nested-anchor}, $S_{S_n}=rn-R_n$ with $0\le R_n\le r(r-1)/2$.
Since $R_n=\sum_{j=1}^{r-1}(r-j)\,d(S_n-j)$, the shift $R_n$ is itself
determined by the local differences $d(S_n-1),\ldots,d(S_n-r+1)$, which
by Lemma~\ref{lem:local_projection}$(ii)$ depend only on the run lengths
$(L_{n-1},\ldots,L_{n-r+2})$.
Set $X:=R_n+i\in[0,D]$.
Then $d(rn+i)=d(S_{S_n}+X)$, and by Lemma~\ref{lem:local_projection}$(i)$ at $A=S_n$,
this is determined by $(L_{S_n},\ldots,L_{S_n+D})$.

\textit{From scale $S_n$ to scale $n$.}
For $0\le k\le D$, Lemma~\ref{lem:runlength-d} gives
$L_{S_n+k}=\sum_{j=0}^{r-1}d(S_n+k-j)$,
so we need $d(S_n+x)$ for $x\in[-(r-1),D]$.
For $x\ge 0$: Lemma~\ref{lem:local_projection}$(i)$ at $A=n$ gives this from
$(L_n,\ldots,L_{n+D})$.
For $x=-1$: $d(S_n-1)=1$.
For $x=-1-Y$ with $1\le Y\le r-2$: Lemma~\ref{lem:local_projection}$(ii)$ at $A=n$
gives this from $(L_{n-1},\ldots,L_{n-Y})$.
All needed $L_m$ have index $m\ge n-r+2\ge 3$ (since $n\ge n_0$), so
Lemma~\ref{lem:runlength-d} applies: $L_m=\sum_{j=0}^{r-1}d(m-j)$,
determined by $d(n-C),\ldots,d(n+D)$.

\textit{Conclusion.}
There exists a Boolean function $\varphi_i:\{0,1\}^{C+D+1}\to\{0,1\}$ such that
\[
  d(rn+i)=\varphi_i\bigl(d(n-C),\ldots,d(n+D)\bigr)\qquad(n\ge n_0).
\]
By Theorem~\ref{thm:AS}, $d$ is $r$-automatic.
Since $a(n)=\sum_{k=0}^{n-1}d(k)$, the sequence $(a(n))$ is $r$-regular
by~\cite[Theorem~2.5]{AS92}.
\end{proof}

\begin{remark}
The proof is qualitative: it produces a finite automaton for $d$ but does not
minimise it. The explicit automata for $r\le 5$ (Sections~\ref{sec:r2}--\ref{sec:r5}
and Appendix~\ref{app:dfao}) are much smaller, which is why the
constructive cases remain informative even after Theorem~\ref{thm:structural_regularity}.
For $r\ge 6$, what remains open is the existence of compact denesting formulas
$a(rn+i)=T_r(n+s_i)+\delta_i+\varepsilon_i(n)$ with few $r$-automatic sequences.
\end{remark}



\section{The case $r=2$}\label{sec:r2}
%-------------------------------------------------------------------

\subsection{Definition and dyadic structure}

Throughout the paper we use the phrase \emph{$r$-adic recurrence} (or
\emph{dyadic} for $r=2$, \emph{triadic} for $r=3$, \emph{quaternary} for $r=4$)
to mean a recurrence that decomposes $n$ according to its base-$r$ representation,
of the form $a(rn+i) = f_i(a(n), a(n\pm 1), \ldots)$. This usage is unrelated to
$p$-adic numbers.

The almost Golomb sequence of order $2$ is defined by $a(0)=0$, $a(1)=1$ and,
for $n\ge 2$, $a(n)$ is the smallest integer $\ge a(n-1)$ such that
\begin{equation}\label{eq:r2def}
  a\bigl(a(n)+a(n-1)\bigr) = n.
\end{equation}
Set $S_n = a(n)+a(n-1)$ for $n\ge 1$, so the defining property reads $a(S_n)=n$.
By Lemma~\ref{lem:anchor-growth}, the indices $S_n$ are pairwise distinct.
This sequence is \oeis{A394217} in the OEIS. The first terms are
\[
  1,2,2,3,4,4,5,6,7,7,8,8,9,10,11,12,13,13,14,14,15,15,16,16,\ldots
\]


\subsection{The explicit dyadic formula}

We derive a closed-form expression for $a(n)$ in each
\emph{dyadic block} $[2^k, 2^{k+1})$, i.e.\ the set of integers
$n$ satisfying $2^k \le n < 2^{k+1}$.

\begin{proposition}\label{prop:dyadic}
  For $k\ge 2$ and $n = 2^k+j$ with $0\le j < 2^k$,
  \[
    a(n) = \begin{cases}
      3\cdot 2^{k-2} + \lceil j/2\rceil & \text{if } 0\le j\le 2^{k-1},\\
      n - 2^{k-1}                       & \text{if } 2^{k-1} < j < 2^k.
    \end{cases}
  \]
\end{proposition}

\begin{proof}
  We prove that the displayed formula defines the almost Golomb sequence of order $2$,
  i.e.\ the earliest monotone sequence satisfying \eqref{eq:r2def}.

  Let $\widetilde a(0)=0$, $\widetilde a(1)=1$, and define $\widetilde a(n)$ for
  $n\ge 2$ by the displayed dyadic formula (with the convention
  $\widetilde a(2)=\widetilde a(3)=2$ when $k=1$).
  We show that $\widetilde a$ satisfies the greedy rule \eqref{eq:r2def}; hence
  $\widetilde a=a$.

  First, $\widetilde a$ is nondecreasing and takes positive integer values,
  by inspection of the formula on each dyadic block.

  Fix $n\ge 2$ and set $T_n := \widetilde a(n)+\widetilde a(n-1)$.
  We claim
  \begin{equation}\label{eq:r2_anchor}
    \widetilde a(T_n)=n
    \qquad\text{and}\qquad
    \widetilde a(T_n-1)=n-1.
  \end{equation}
  This will imply both that $\widetilde a$ satisfies the implicit equation
  $\widetilde a(\widetilde a(n)+\widetilde a(n-1))=n$ and that the choice
  $\widetilde a(n)$ is minimal among all $m\ge \widetilde a(n-1)$ achieving it.

  Write $n=2^k+j$ with $k\ge 2$ and $0\le j<2^k$.
  We treat three cases.

  \textit{Case 1: $j=0$.} Then $\widetilde a(n)=3\cdot 2^{k-2}$ and
  $n-1=2^k-1$ lies in the second half of the previous dyadic block, so
  $\widetilde a(n-1)=(n-1)-2^{k-2}=3\cdot 2^{k-2}-1$,
  hence $T_n=3\cdot 2^{k-1}-1$.
  Now $T_n=2^k+(2^{k-1}-1)$ lies in the first half of $[2^k,2^{k+1})$, so
  \[
    \widetilde a(T_n)=3\cdot 2^{k-2}+\Bigl\lceil\frac{2^{k-1}-1}{2}\Bigr\rceil
    =3\cdot 2^{k-2}+2^{k-2}=2^k=n.
  \]
  Moreover $T_n-1=2^k+(2^{k-1}-2)$ is still in the first half, hence
  \[
    \widetilde a(T_n-1)=3\cdot 2^{k-2}+\Bigl\lceil\frac{2^{k-1}-2}{2}\Bigr\rceil
    =3\cdot 2^{k-2}+(2^{k-2}-1)=2^k-1=n-1.
  \]

  \textit{Case 2: $1\le j\le 2^{k-1}$.} Both $n$ and $n-1$ lie in the first half
  of the same dyadic block, so
  $\widetilde a(n)=3\cdot 2^{k-2}+\lceil j/2\rceil$ and
  $\widetilde a(n-1)=3\cdot 2^{k-2}+\lceil(j-1)/2\rceil$.
  Using $\lceil j/2\rceil+\lceil(j-1)/2\rceil=j$, we get $T_n=3\cdot 2^{k-1}+j$.
  If $j<2^{k-1}$, then $T_n=2^k+(2^{k-1}+j)$ lies in the second half of
  $[2^k,2^{k+1})$, giving $\widetilde a(T_n)=T_n-2^{k-1}=2^k+j=n$ and
  $\widetilde a(T_n-1)=n-1$.
  If $j=2^{k-1}$, then $T_n=2^{k+1}$ and
  $\widetilde a(T_n)=\widetilde a(2^{k+1})=3\cdot 2^{k-1}=2^k+2^{k-1}=n$;
  while $T_n-1=2^{k+1}-1$ lies in the second half of $[2^k,2^{k+1})$ and
  $\widetilde a(T_n-1)=(2^{k+1}-1)-2^k=2^k-1=n-1$.

  \textit{Case 3: $2^{k-1}<j<2^k$.} Write $j=2^{k-1}+u$ with $1\le u\le 2^{k-1}-1$.
  Both $n$ and $n-1$ lie in the second half of $[2^k,2^{k+1})$, so
  $\widetilde a(n)=n-2^{k-1}$ and $\widetilde a(n-1)=(n-1)-2^{k-1}$.
  Hence $T_n=2n-1-2^k=2^{k+1}+(2u-1)$.
  Since $1\le 2u-1\le 2^k-3$, this lies in the first half of $[2^{k+1},2^{k+2})$, and
  \[
    \widetilde a(T_n)=3\cdot 2^{k-1}+\Bigl\lceil\frac{2u-1}{2}\Bigr\rceil
    =3\cdot 2^{k-1}+u=2^k+j=n,
  \]
  and similarly $\widetilde a(T_n-1)=3\cdot 2^{k-1}+(u-1)=n-1$.

  This proves \eqref{eq:r2_anchor} for all $n\ge 2$.
  Hence $\widetilde a(\widetilde a(n)+\widetilde a(n-1))=n$.
  Moreover, for any $m$ with $\widetilde a(n-1)\le m\le \widetilde a(n)-1$,
  we have $m+\widetilde a(n-1)\le T_n-1$, hence
  $\widetilde a(m+\widetilde a(n-1))\le \widetilde a(T_n-1)=n-1$,
  so $m$ cannot satisfy \eqref{eq:r2def}.
  Therefore $\widetilde a(n)$ is the smallest integer $\ge \widetilde a(n-1)$
  satisfying \eqref{eq:r2def}, which is exactly the greedy rule.
  Hence $\widetilde a=a$.
\end{proof}

Specialising to the endpoints of dyadic blocks gives two exact values.

\begin{corollary}\label{cor:pivots}
  For all $k\ge 2$: $a(2^k) = 3\cdot 2^{k-2}$ and $a(3\cdot 2^{k-1}) = 2^k$.
\end{corollary}

\subsection{The denesting theorem}

The closed form yields the denesting formulas by a direct comparison.

\begin{theorem}\label{thm:r2}
  The almost Golomb sequence of order $2$ satisfies
  \begin{align}
    a(2n)   &= a(n)+a(n+1)-1, \qquad (n\ge 1), \label{eq:r2even}\\
    a(2n+1) &= a(n)+a(n+1), \qquad (n\ge 2).   \label{eq:r2odd}
  \end{align}
  In particular $(a(n))$ is $2$-regular in the sense of Allouche--Shallit~\cite{AS92}.
\end{theorem}

\begin{proof}
  We use Proposition~\ref{prop:dyadic} throughout. Write $n=2^k+j$ with $k\ge 2$
  and $0\le j<2^k$. We treat four cases according to the position of $j$ in the
  dyadic block. The key identity used throughout is
  \begin{equation}\label{eq:ceil-cancel}
    \Bigl\lceil\frac{j}{2}\Bigr\rceil+\Bigl\lceil\frac{j+1}{2}\Bigr\rceil = j+1,
  \end{equation}
  valid for all integers $j\ge 0$.


  \textit{Case 1: $0\le j\le 2^{k-1}-1$ (strict first half).}
  Both $j$ and $j+1$ lie in $[0,2^{k-1}]$, so
  \[
    a(n)=3\cdot 2^{k-2}+\Bigl\lceil\tfrac{j}{2}\Bigr\rceil,\qquad
    a(n+1)=3\cdot 2^{k-2}+\Bigl\lceil\tfrac{j+1}{2}\Bigr\rceil.
  \]
  By \eqref{eq:ceil-cancel}: $a(n)+a(n+1)-1=3\cdot 2^{k-1}+j$ and
  $a(n)+a(n+1)=3\cdot 2^{k-1}+j+1$.
  Since $2j\in[0,2^k-2]$ (first half of block $k+1$):
  $a(2n)=3\cdot 2^{k-1}+j$. Since $2j+1\in[1,2^k-1]$ (first half):
  $a(2n+1)=3\cdot 2^{k-1}+j+1$.


  \textit{Case 2: $j=2^{k-1}$ (right endpoint of first half).}
  $a(n)=3\cdot 2^{k-2}+2^{k-2}=2^k$. Since $j+1=2^{k-1}+1>2^{k-1}$
  (second half): $a(n+1)=n+1-2^{k-1}=2^k+1$.
  Hence $a(n)+a(n+1)-1=2^{k+1}$ and $a(n)+a(n+1)=2^{k+1}+1$.
  Since $2j=2^k=2^{(k+1)-1}$ (boundary of first half in block $k+1$):
  $a(2n)=3\cdot 2^{k-1}+2^{k-1}=2^{k+1}$.
  Since $2j+1=2^k+1>2^k$ (second half): $a(2n+1)=2n+1-2^k=2^{k+1}+1$.


  \textit{Case 3: $2^{k-1}<j<2^k-1$ (strict second half).}
  Both $j,j+1>2^{k-1}$, so $a(n)=n-2^{k-1}$ and $a(n+1)=n+1-2^{k-1}$.
  Hence $a(n)+a(n+1)-1=2n-2^k$ and $a(n)+a(n+1)=2n-2^k+1$.
  Since $2j>2^k$ (second half of block $k+1$): $a(2n)=2n-2^k$
  and $a(2n+1)=2n+1-2^k$.


  \textit{Case 4: $j=2^k-1$ (last element of block $k$).}
  $a(n)=n-2^{k-1}=3\cdot 2^{k-1}-1$. Now $n+1=2^{k+1}$ starts block $k+1$
  with $j'=0$, giving $a(n+1)=3\cdot 2^{k-1}$.
  Hence $a(n)+a(n+1)-1=3\cdot 2^k-2$ and $a(n)+a(n+1)=3\cdot 2^k-1$.
  Since $2j=2^{k+1}-2>2^k$ (second half): $a(2n)=2n-2^k=3\cdot 2^k-2$
  and $a(2n+1)=2n+1-2^k=3\cdot 2^k-1$.


  The special case $k=1$ (i.e.\ $n\in\{2,3\}$) is verified directly:
  $a(4)=a(2)+a(3)-1=3$, $a(5)=a(2)+a(3)=4$,
  $a(6)=a(3)+a(4)-1=4$, $a(7)=a(3)+a(4)=5$;
  these match the sequence.
\end{proof}

\begin{remark}
  Formulas \eqref{eq:r2even}--\eqref{eq:r2odd} show that every subsequence in the
  $2$-kernel of $(a(n))$ (see the discussion preceding Theorem~\ref{thm:AS} for the terminology)
  can be expressed as a $\mathbb{Z}$-linear combination of
  $(a(n))$, $(a(n+1))$, and the constant sequence $1$ (the $-1$ in \eqref{eq:r2even}
  necessitates this constant generator). Hence $(a(n))$ is $2$-regular.
  Since $a(n)$ is unbounded it is not $2$-automatic.
\end{remark}

\subsection{A Mallows-type nested recurrence}

Golomb's sequence satisfies the Mallows recurrence
$G(n+1)=1+G(n+1-G(G(n)))$. The almost Golomb sequence of order~$2$
admits an analogous nested recurrence.

\begin{proposition}\label{prop:mallows-r2}
For every $n\ge 4$,
\begin{equation}\label{eq:r2-mallows}
a(n+1)
\;=\;
1+a\!\Bigl(n+1-a\bigl(a(n)+1\bigr)+a\bigl(a(n)-1\bigr)\Bigr).
\end{equation}
\end{proposition}

\begin{proof}
Set $m:=a(n)$ and let $N_2(m):=a(m+1)-a(m-1)$ be the multiplicity of
the value~$m$, i.e.\ the length of the run of~$m$
(Lemma~\ref{lem:runlength-d} with $r=2$).
By Theorem~\ref{thm:global_run}, the differences $d\in\{0,1\}$ give
$N_2(m)=d(m-2)+d(m-1)\in\{1,2\}$.
It suffices to show
\begin{equation}\label{eq:r2-mallows-aux}
a(n+1)=1+a\bigl(n+1-N_2(m)\bigr),
\end{equation}
since substituting $N_2(m)=a(m+1)-a(m-1)$ with $m=a(n)$
yields~\eqref{eq:r2-mallows}.

If $N_2(m)=1$: the value $m$ occurs once, at position~$n$,
so $a(n+1)=m+1$ and $1+a(n)=1+m=m+1$.

If $N_2(m)=2$: the value $m$ occupies two consecutive positions.
If $n$ is the first, then $a(n-1)=m-1$ and $a(n+1)=m$,
giving $1+a(n-1)=m$.
If $n$ is the second, then $a(n-1)=m$ and $a(n+1)=m+1$,
giving $1+a(n-1)=m+1$.
\end{proof}

\begin{remark}
Since $a(n)<n$ for $n\ge 3$, the right-hand side of~\eqref{eq:r2-mallows}
involves only values $a(k)$ with $k\le n$.
Hence~\eqref{eq:r2-mallows} together with the initial conditions
$a(1)=1$, $a(2)=2$, $a(3)=2$, $a(4)=3$ provides an alternative
definition of the almost Golomb sequence of order $2$,
without reference to the greedy rule or the implicit
equation~\eqref{eq:r2def}.
\end{remark}

\section{The case $r=3$}\label{sec:r3}
%-------------------------------------------------------------------

\subsection{Definition and the correction sequence}

The almost Golomb sequence of order $3$ is defined by $a(1)=1$, $a(k)=0$ for $k<1$,
and for $n\ge 2$, $a(n)$ is the smallest integer $\ge a(n-1)$ such that
\begin{equation}\label{eq:r3def}
  a\bigl(a(n)+a(n-1)+a(n-2)\bigr)=n.
\end{equation}
Set $S_n:=a(n)+a(n-1)+a(n-2)$, so that $a(S_n)=n$ for all $n\ge 1$.
This sequence is \oeis{A394218} in the OEIS. It begins
$1,2,2,2,3,4,5,5,6,6,6,7,7,8,8,9,10,11,12,13,13,14,15,15,\ldots$

Define
\begin{equation}\label{eq:Ik}
  I_k=\left[\frac{11\cdot 3^k-1}{2},\,\frac{13\cdot 3^k-3}{2}\right]\cap\mathbb{Z},
  \qquad
  \mathcal{I}=\bigcup_{k\ge 0} I_k,
  \qquad
  \varepsilon(n)=\mathbf{1}_{\mathcal{I}}(n).
\end{equation}
Thus $I_0=\{5\}$ and $|I_k|=3^k$ for every $k\ge 0$.

\subsection{Automaticity of the correction sequence}

\begin{lemma}\label{lem:selfsim}
For every $k\ge 0$, $I_{k+1}=(3I_k+1)\cup(3I_k+2)\cup(3I_k+3)$.
\end{lemma}

\begin{proof}
With $\ell_k=(11\cdot 3^k-1)/2$ and $u_k=(13\cdot 3^k-3)/2$,
one has $\ell_{k+1}=3\ell_k+1$ and $u_{k+1}=3u_k+3$, giving
$I_{k+1}=[3\ell_k+1,3u_k+3]\cap\mathbb{Z}=(3I_k+1)\cup(3I_k+2)\cup(3I_k+3)$.
\end{proof}

\begin{proposition}\label{prop:eps-auto}
The sequence $\varepsilon$ satisfies $\varepsilon(n)=0$ for $n<5$, $\varepsilon(5)=1$,
and for every $m\ge 2$:
\begin{equation}\label{eq:eps-rec}
  \varepsilon(3m)=\varepsilon(m-1),\qquad
  \varepsilon(3m+1)=\varepsilon(m),\qquad
  \varepsilon(3m+2)=\varepsilon(m).
\end{equation}
In particular, $\varepsilon$ is $3$-automatic.
\end{proposition}

\begin{proof}
From Lemma~\ref{lem:selfsim}: $3m\in 3I_k+3\iff m-1\in I_k$, and
$3m+1\in 3I_k+1$, $3m+2\in 3I_k+2$ both reduce to $m\in I_k$.
Since $\varepsilon$ is binary and obeys the digit-by-digit recurrences
\eqref{eq:eps-rec}, it is generated by a finite automaton reading the base-$3$
expansion of $n$.
\end{proof}

\begin{proposition}\label{prop:eps-dfao}
The sequence $\varepsilon(n)=\mathbf{1}_{\mathcal{I}}(n)$ is generated by the following
DFAO, reading the base-$3$ expansion of $n$ from most significant to least significant digit.
Let $Q=\{q_0,q_1,q_2,c_{=},c_{>},c_{<},e_{=},e_{>},e_{<},r\}$, $\Sigma=\{0,1,2\}$,
initial state $q_0$, output map $\tau(c_{=})=\tau(c_{>})=\tau(e_{<})=1$ and $\tau(q)=0$
otherwise. The transition table is:
\[
\begin{array}{c|ccc}
      & 0 & 1 & 2\\ \hline
q_0   & r & q_1 & q_2\\
q_1   & r & r   & c_{=}\\
q_2   & e_{=} & r & r\\
c_{=} & c_{<} & c_{=} & c_{>}\\
c_{<} & c_{<} & c_{<} & c_{<}\\
c_{>} & c_{>} & c_{>} & c_{>}\\
e_{=} & e_{<} & e_{=} & e_{>}\\
e_{<} & e_{<} & e_{<} & e_{<}\\
e_{>} & e_{>} & e_{>} & e_{>}\\
r     & r & r & r
\end{array}
\]
For every $n\ge 1$, the automaton outputs $\varepsilon(n)=\mathbf{1}_{\mathcal{I}}(n)$.
\end{proposition}

\begin{proof}
The bounds $\ell_k=(11\cdot 3^k-1)/2=(12\underbrace{1\cdots1}_{k})_3$ and
$u_k=(13\cdot 3^k-3)/2=(20\underbrace{1\cdots1}_{k-1}0)_3$ show that
$n\in I_k$ iff either (i)~$\operatorname{rep}_3(n)=12u$ with $|u|=k$ and $u\ge_{\rm lex}1^k$,
or (ii)~$\operatorname{rep}_3(n)=20u$ with $|u|=k$ and $u<_{\rm lex}1^k$.
States $q_1,q_2$ detect prefixes $12,20$; states $c_{=},c_{>},c_{<}$ (resp.\ $e_{=},e_{>},e_{<}$)
track the lexicographic comparison of the suffix with $1^k$.
Accepting states $c_{=},c_{>}$ cover case~(i); $e_{<}$ covers case~(ii).
\end{proof}

\subsection{A strong gap lemma}

\begin{lemma}\label{lem:r3-gap}
For every $k\ge 2$, $\varepsilon(n)=0$ for all $3^k-1\le n\le 5\cdot 3^{k-1}$.
\end{lemma}

\begin{proof}
By induction on $k$. The base case $k=2$ is $\varepsilon(8)=\cdots=\varepsilon(15)=0$,
verified directly. Assuming the result at rank $k$, let $3^{k+1}-1\le n\le 5\cdot 3^k$
and write $n=3m+r$.
If $r\in\{1,2\}$: $\varepsilon(n)=\varepsilon(m)$ and $3^k-1\le m\le 5\cdot 3^{k-1}$.
If $r=0$: $\varepsilon(n)=\varepsilon(m-1)$ and $3^k-1\le m-1\le 5\cdot 3^{k-1}-1$.
In all cases the hypothesis gives $\varepsilon=0$.
\end{proof}

\subsection{Explicit formulas on the first two thirds of each block}

\begin{lemma}\label{lem:r3-good-zone}
For every $k\ge 2$ and every $n$ with $3^k-1\le n\le 5\cdot 3^{k-1}$,
the almost Golomb sequence satisfies
\begin{align}
  a(3n)   &= a(n-2)+a(n-1)+a(n)+1, \label{eq:r3-good0}\\
  a(3n+1) &= a(n-1)+a(n)+a(n+1),   \label{eq:r3-good1}\\
  a(3n+2) &= a(n)+a(n+1)+a(n+2)-1. \label{eq:r3-good2}
\end{align}
\end{lemma}

\begin{proof}
This is the specialisation of Theorem~\ref{thm:r3}
(whose proof does not depend on the present lemma) to the zone
where $\varepsilon\equiv 0$: Lemma~\ref{lem:r3-gap} gives
$\varepsilon(n-1)=\varepsilon(n)=0$ for $3^k-1\le n\le 5\cdot 3^{k-1}$,
so \eqref{eq:r3a}--\eqref{eq:r3c} reduce to
\eqref{eq:r3-good0}--\eqref{eq:r3-good2}.
\end{proof}

\begin{proposition}\label{prop:r3-block}
For every $k\ge 2$, set $A_k:=a(3^k)$. Then $A_2=6$, $A_{k+1}=3A_k-1$, and
\[
  A_k=\frac{11\cdot 3^{k-2}+1}{2}.
\]
Moreover $a(3^k-2)=a(3^k-1)=A_k-1$, and for $0\le j\le 2\cdot 3^{k-1}$:
\begin{equation}\label{eq:r3-block}
a(3^k+j)=
\begin{cases}
  A_k+\left\lfloor \dfrac{j}{3}\right\rfloor,
    & 0\le j\le 3^{k-1}-1,\\[2ex]
  A_k+3^{k-2}+\left\lfloor \dfrac{j-3^{k-1}}{2}\right\rfloor,
    & 3^{k-1}\le j\le 2\cdot 3^{k-1}.
\end{cases}
\end{equation}
\end{proposition}

\begin{proof}
By induction on $k$. The base case $k=2$ is verified from
the values $a(9)=\cdots=a(11)=6$, $a(12)=a(13)=7$, $a(14)=a(15)=8$.

Assume \eqref{eq:r3-block} at rank $k$.
By Lemma~\ref{lem:r3-good-zone}, for every $n\in[3^k-1,\,5\cdot 3^{k-1}]$,
\begin{align}
  a(3n)   &= a(n-2)+a(n-1)+a(n)+1, \\
  a(3n+1) &= a(n-1)+a(n)+a(n+1),   \\
  a(3n+2) &= a(n)+a(n+1)+a(n+2)-1.
\end{align}
Left boundary: at $n=3^k-1$, \eqref{eq:r3-good1}--\eqref{eq:r3-good2} give
$a(3^{k+1}-2)=a(3^{k+1}-1)=3A_k-2=A_{k+1}-1$, and \eqref{eq:r3-good0} gives
$a(3^{k+1})=3A_k-1=A_{k+1}$.

First subinterval: for $q=3^k+v$, $0\le v<3^{k-1}$, the floor sum
$\lfloor(v-2)/3\rfloor+\lfloor(v-1)/3\rfloor+\lfloor v/3\rfloor=v-2$
and its two analogues give $a(3q)=a(3q+1)=a(3q+2)=A_{k+1}+v$.

Second subinterval: for $q=4\cdot 3^{k-1}+v$, $0\le v\le 3^{k-1}$, the analogous
floor sums with denominator $2$ give
$a(4\cdot 3^k+s)=A_{k+1}+3^{k-1}+\lfloor s/2\rfloor$ for $0\le s\le 3^k$.

Solving $A_{k+1}=3A_k-1$ with $A_2=6$ gives $A_k=(11\cdot 3^{k-2}+1)/2$.
\end{proof}

\begin{corollary}\label{cor:r3-special}
For every $k\ge 2$:
\[
  a(3^k)=\frac{11\cdot 3^{k-2}+1}{2},
  \qquad
  a(5\cdot 3^{k-1})=8\cdot 3^{k-2},
  \qquad
  a(5\cdot 3^{k-1}-1)=8\cdot 3^{k-2}-1.
\]
\end{corollary}

\begin{proof}
Take $j=2\cdot 3^{k-1}$ and $j=2\cdot 3^{k-1}-1$ in \eqref{eq:r3-block}.
Since $3^{k-1}$ is odd, $\lfloor 3^{k-1}/2\rfloor=(3^{k-1}-1)/2$.
\end{proof}

\subsection{Triadic denesting}

\begin{lemma}\label{lem:subdiag}
For every $r\ge 2$ and every $n\ge 3$, $a(n)\le n-1$.
\end{lemma}

\begin{proof}
Since $a(3)=2$ and $a(n+1)-a(n)\in\{0,1\}$ by
Theorem~\ref{thm:global_run}, an immediate induction gives
$a(n)\le n-1$ for all $n\ge 3$.
\end{proof}

\begin{theorem}\label{thm:r3}
For every $n\ge 2$:
\begin{align}
  a(3n)   &= a(n-2)+a(n-1)+a(n)+1+\varepsilon(n-1), \label{eq:r3a}\\
  a(3n+1) &= a(n-1)+a(n)+a(n+1),                    \label{eq:r3b}\\
  a(3n+2) &= a(n)+a(n+1)+a(n+2)-1-\varepsilon(n).   \label{eq:r3c}
\end{align}
\end{theorem}

\begin{proof}
We prove the three identities \eqref{eq:r3a}--\eqref{eq:r3c} together with
the \emph{run-length invariant}: defining
\begin{equation}\label{eq:Ldef}
  L_n\;:=\;S_{n+1}-S_n\;=\;a(n+1)-a(n-2)\;=\;d(n-2)+d(n-1)+d(n)
  \;\in\;\{1,2,3\}
\end{equation}
(the length of the run of value~$n$, for $n\ge 3$), we prove that for every $n\ge 3$,
\begin{equation}\label{eq:r3-Linv}
  L_n \le 2+\varepsilon(n-1),
  \qquad\text{with equality } L_n=3 \text{ whenever } \varepsilon(n-1)=1.
\end{equation}
The proof is by strong induction on~$n$.
The hypothesis $\mathcal{H}_n$ asserts that the three denesting formulas hold
at all ranks $m<n$, and that the invariant~\eqref{eq:r3-Linv} holds for
all $m$ with $3\le m\le n+1$.
(The invariant at $m\le n+1$ involves only $a(k)$ for $k\le m+1\le n+2<3n$,
hence values already determined at a previous induction step.)

\textit{Base cases.}
Direct computation for $2\le n\le 30$ verifies the denesting formulas
and the invariant~\eqref{eq:r3-Linv}.
This threshold covers all $n$ up to well past $I_1=[16,18]$,
the first non-trivial component of $\mathcal{I}$. For $n\ge 31$, every index
encountered in the induction step satisfies the hypotheses of
Proposition~\ref{prop:eps-auto} (valid for $m\ge 2$) and
Lemma~\ref{lem:r3-gap} (valid for $k\ge 2$).

\textit{Induction step.}
Fix $n\ge 31$ and assume $\mathcal{H}_n$.
We determine the three new values $a(3n)$, $a(3n+1)$, $a(3n+2)$
via Lemma~\ref{lem:anchor-characterization}: for $M\ge 3$,
\begin{equation}\label{eq:anchor-char}
  a(K)=M \iff S_M\le K<S_{M+1}.
\end{equation}
For each candidate $M$ below, one has $M<3n$.
Indeed, by Lemma~\ref{lem:subdiag} ($a(t)\le t-1$ for $t\ge 3$):
\[
  M_1 = S_{n+1} \le n+(n{-}1)+(n{-}2)=3n{-}3,
\]
\[
  M_0 = S_n+1+\varepsilon(n{-}1) \le (n{-}1)+(n{-}2)+(n{-}3)+2 = 3n{-}4,
\]
\[
  M_2 = S_{n+2}-1-\varepsilon(n) \le (n{+}1)+n+(n{-}1)-1 = 3n{-}1.
\]
Hence the values $a(M),a(M-1),a(M-2)$ entering $S_M$ are
already determined by the induction hypothesis.

\textit{Case $a(3n+1)$.}
The candidate is $M_1:=a(n-1)+a(n)+a(n+1)=S_{n+1}$.
By the run structure, $a(S_{n+1})=n+1$ and $a(S_{n+1}-1)=n$.
If $L_n\ge 2$ then $a(S_{n+1}-2)=n$; if $L_n=1$ then
$S_{n+1}-2=S_n-1$ and $a(S_n-1)=n-1$.
Hence $S_{M_1}\in\{3n,\,3n+1\}$, so $S_{M_1}\le 3n+1$.
For the upper bound,
$L_{M_1}=a(S_{n+1}+1)-a(S_{n+1}-2)$.
If $L_{n+1}\ge 2$ then $a(S_{n+1}+1)=n+1$;
if $L_{n+1}=1$ then $a(S_{n+1}+1)=n+2$.
Checking all combinations: $S_{M_1+1}=S_{M_1}+L_{M_1}\ge 3n+2>3n+1$.
By~\eqref{eq:anchor-char}, $a(3n+1)=S_{n+1}$,
which is~\eqref{eq:r3b}.

\textit{Case $a(3n)$.}
The candidate is $M_0:=a(n-2)+a(n-1)+a(n)+1+\varepsilon(n-1)=S_n+1+\varepsilon(n-1)$.

\textit{Subcase $\varepsilon(n-1)=0$.}
Then $M_0=S_n+1$ and $L_n\le 2$ by~\eqref{eq:r3-Linv}.
Since $a(S_n)=n$ and $a(S_n-1)=n-1$:
if $L_n=2$ then $a(M_0)=n$, giving
$S_{M_0}=n+n+(n-1)=3n-1$;
if $L_n=1$ then $M_0=S_{n+1}$, $a(M_0)=n+1$, giving
$S_{M_0}=(n+1)+n+(n-1)=3n$.
In both cases $S_{M_0}\le 3n$.
For the upper bound: when $L_n=2$,
$a(M_0+1)=a(S_{n+1})=n+1$, so $L_{M_0}=(n+1)-(n-1)=2$
and $S_{M_0+1}=3n+1$;
when $L_n=1$,
$L_{M_0}=a(S_{n+1}+1)-a(S_n-1)\ge(n+1)-(n-1)=2$
and $S_{M_0+1}\ge 3n+2$.
In both cases $S_{M_0+1}>3n$.

\textit{Subcase $\varepsilon(n-1)=1$.}
Then $M_0=S_n+2$ and $L_n=3$ by~\eqref{eq:r3-Linv}.
All three positions $S_n,\,S_n+1,\,S_n+2$ carry value~$n$, so
$S_{M_0}=n+n+n=3n$.
Moreover $M_0+1=S_n+3=S_{n+1}$, hence $a(M_0+1)=n+1$,
giving $L_{M_0}=(n+1)-n=1$ and $S_{M_0+1}=3n+1$.

In all subcases $S_{M_0}\le 3n<S_{M_0+1}$,
proving $a(3n)=M_0$, which is~\eqref{eq:r3a}.

\textit{Case $a(3n+2)$.}
The candidate is $M_2:=a(n)+a(n+1)+a(n+2)-1-\varepsilon(n)=S_{n+2}-1-\varepsilon(n)$.

\textit{Subcase $\varepsilon(n)=0$.}
Then $M_2=S_{n+2}-1$ and $L_{n+1}\le 2$.
Here $a(M_2)=a(S_{n+2}-1)=n+1$ (last position of the run of $n+1$).
If $L_{n+1}=2$: $a(M_2-1)=n+1$, $a(M_2-2)=a(S_{n+1}-1)=n$,
giving $S_{M_2}=3n+2$.
If $L_{n+1}=1$: $a(M_2-1)=a(S_{n+1}-1)=n$.
\quad If $L_n\ge 2$: $a(M_2-2)=n$, giving $S_{M_2}=3n+1$.
\quad If $L_n=1$: $a(M_2-2)=a(S_n-1)=n-1$, giving $S_{M_2}=3n$.
In all cases $S_{M_2}\le 3n+2$.
For the upper bound, $M_2+1=S_{n+2}$, so
$S_{M_2+1}=S_{S_{n+2}}$.
By Lemma~\ref{lem:nested-anchor},
$S_{S_{n+2}}=3(n+2)-2d(S_{n+2}-1)-d(S_{n+2}-2)=3n+4-d(S_{n+2}-2)$
(using $d(S_{n+2}-1)=1$, Lemma~\ref{lem:local_projection}$(ii)$).
By Lemma~\ref{lem:local_projection}$(ii)$, $d(S_{n+2}-2)=1$
iff $L_{n+1}=1$.
Hence $S_{M_2+1}\ge 3n+3>3n+2$.

\textit{Subcase $\varepsilon(n)=1$.}
Then $M_2=S_{n+2}-2$ and $L_{n+1}=3$ by~\eqref{eq:r3-Linv}.
The run of $n+1$ occupies $[S_{n+1},\,S_{n+1}+2]$, so
$a(M_2)=a(S_{n+1}+1)=n+1$,
$a(M_2-1)=a(S_{n+1})=n+1$,
$a(M_2-2)=a(S_{n+1}-1)=n$,
giving $S_{M_2}=(n+1)+(n+1)+n=3n+2$.
For the upper bound,
$L_{M_2}=a(S_{n+2}-1)-a(S_{n+1}-2)=(n+1)-a(S_{n+1}-2)$.
If $L_n\ge 2$: $a(S_{n+1}-2)=n$, $L_{M_2}=1$;
if $L_n=1$: $a(S_{n+1}-2)=a(S_n-1)=n-1$, $L_{M_2}=2$.
In both cases $S_{M_2+1}=S_{M_2}+L_{M_2}\ge 3n+3>3n+2$.

In all subcases $S_{M_2}\le 3n+2<S_{M_2+1}$,
proving $a(3n+2)=M_2$, which is~\eqref{eq:r3c}.

\textit{Propagation of the run-length invariant.}
With the denesting established at rank~$n$, a direct calculation
using \eqref{eq:r3a}--\eqref{eq:r3c} at ranks $n-1$ and $n$ gives:
\begin{align}
  L_{3n-1} &= L_{n-1}+\varepsilon(n-1)-\varepsilon(n-2),\label{eq:Lprop1}\\
  L_{3n}   &= L_n,                                       \label{eq:Lprop2}\\
  L_{3n+1} &= L_{n+1}+\varepsilon(n-1)-\varepsilon(n).   \label{eq:Lprop3}
\end{align}
The $\varepsilon$-recurrence~\eqref{eq:eps-rec} gives
$\varepsilon(3n-2)=\varepsilon(3n-1)=\varepsilon(3n)=\varepsilon(n-1)$.

\emph{Upper bounds.}
$L_{3n-1}\le(2+\varepsilon(n-2))+\varepsilon(n-1)-\varepsilon(n-2)=2+\varepsilon(n-1)=2+\varepsilon(3n-2)$.
$L_{3n}=L_n\le 2+\varepsilon(n-1)=2+\varepsilon(3n-1)$.
$L_{3n+1}\le(2+\varepsilon(n))+\varepsilon(n-1)-\varepsilon(n)=2+\varepsilon(n-1)=2+\varepsilon(3n)$.

\emph{Exactness when $\varepsilon=1$.}
Suppose $\varepsilon(3n-2)=\varepsilon(n-1)=1$.
By~\eqref{eq:r3-Linv} at~$n$, $L_n=3$, hence $d(n-2)=d(n-1)=d(n)=1$.
Then $L_{n-1}=d(n-3)+d(n-2)+d(n-1)\ge 2$.
Combined with $L_{n-1}\le 2+\varepsilon(n-2)$:
if $\varepsilon(n-2)=1$ then $L_{n-1}=3$ and $L_{3n-1}=3+1-1=3$;
if $\varepsilon(n-2)=0$ then $L_{n-1}=2$ and $L_{3n-1}=2+1-0=3$.

Suppose $\varepsilon(3n-1)=\varepsilon(n-1)=1$.
Then $L_{3n}=L_n=3$ directly.

Suppose $\varepsilon(3n)=\varepsilon(n-1)=1$.
By~\eqref{eq:r3-Linv} at~$n$, $L_n=3$, hence $d(n-1)=d(n)=1$.
Then $L_{n+1}=d(n-1)+d(n)+d(n+1)\ge 2$.
Combined with $L_{n+1}\le 2+\varepsilon(n)$:
if $\varepsilon(n)=1$ then $L_{n+1}=3$ and $L_{3n+1}=3+1-1=3$;
if $\varepsilon(n)=0$ then $L_{n+1}=2$ and $L_{3n+1}=2+1-0=3$.

This closes the strong induction.
\end{proof}

\subsection{$3$-regularity}

\begin{corollary}\label{cor:r3regular}
The almost Golomb sequence of order $3$ is $3$-regular.
\end{corollary}

\begin{proof}
Since modifying finitely many terms does not affect $3$-regularity,
we work on the range where Theorem~\ref{thm:r3} applies.
Define $A_c(n):=a(n+c)$ for $c\in\{-2,-1,0,1,2\}$,
$E_-(n):=\varepsilon(n-1)$, $E_0(n):=\varepsilon(n)$, and $\mathbf{1}(n):=1$.
Let $M:=\langle A_{-2},A_{-1},A_0,A_1,A_2,E_-,E_0,\mathbf{1}\rangle_{\mathbb{Z}}$
and $(T_rf)(n):=f(3n+r)$.

From \eqref{eq:eps-rec}:
$T_0E_-=E_-$, $T_1E_-=E_-$, $T_2E_-=E_0$,
$T_0E_0=E_-$, $T_1E_0=E_0$, $T_2E_0=E_0$,
$T_r\mathbf{1}=\mathbf{1}$.

From \eqref{eq:r3a}--\eqref{eq:r3c} (shifting by $c$):
\begin{align*}
T_0A_{-2}&=A_{-2}+A_{-1}+A_0, &
T_1A_{-2}&=A_{-1}+A_0+A_1-\mathbf{1}-E_-, \\
T_2A_{-2}&=A_{-2}+A_{-1}+A_0+\mathbf{1}+E_-, &
T_0A_{-1}&=A_{-1}+A_0+A_1-\mathbf{1}-E_-, \\
T_1A_{-1}&=A_{-2}+A_{-1}+A_0+\mathbf{1}+E_-, &
T_2A_{-1}&=A_{-1}+A_0+A_1, \\
T_0A_0   &=A_{-2}+A_{-1}+A_0+\mathbf{1}+E_-, &
T_1A_0   &=A_{-1}+A_0+A_1, \\
T_2A_0   &=A_0+A_1+A_2-\mathbf{1}-E_0, &
T_0A_1   &=A_{-1}+A_0+A_1, \\
T_1A_1   &=A_0+A_1+A_2-\mathbf{1}-E_0, &
T_2A_1   &=A_{-1}+A_0+A_1+\mathbf{1}+E_0, \\
T_0A_2   &=A_0+A_1+A_2-\mathbf{1}-E_0, &
T_1A_2   &=A_{-1}+A_0+A_1+\mathbf{1}+E_0, \\
T_2A_2   &=A_0+A_1+A_2.
\end{align*}
Every image lies in $M$, so $T_r(M)\subseteq M$ for $r=0,1,2$.
Every element of the $3$-kernel of $a$ has the form $T_{r_1}\cdots T_{r_k}A_0\in M$.
Since $M$ is finitely generated, the $3$-kernel spans a finite-rank module,
so $(a(n))$ is $3$-regular.
\end{proof}

\begin{remark}
The ratio $a(n)/n$ does not converge for $r=3$: one has
$a(5\cdot 3^k)/(5\cdot 3^k)=8/15$ and $a(8\cdot 3^k)/(8\cdot 3^k)=5/8$
for all $k\ge 1$. This is established in Section~\ref{sec:nonconv}
(Theorem~\ref{thm:ratio3}) via a propagation argument using
Corollary~\ref{cor:r3-special}.
\end{remark}


\section{The case $r=4$}\label{sec:r4}
%-------------------------------------------------------------------

\subsection{The denesting theorem}

The almost Golomb sequence of order $4$ is defined by $a(1)=1$, $a(k)=0$ for $k<1$,
and for $n\ge 2$, $a(n)$ is the smallest integer $\ge a(n-1)$ such that
\begin{equation}\label{eq:r4def}
  a\bigl(a(n)+a(n-1)+a(n-2)+a(n-3)\bigr) = n.
\end{equation}
This sequence is \oeis{A394219} in the OEIS. The first terms are
\[
  1,2,2,2,3,3,4,4,5,6,6,7,7,8,8,9,9,9,10,10,11,11,11,12,\ldots
\]

As in the triadic case, the denesting of \eqref{eq:r4def} produces four formulas,
each with a binary correction sequence.
In the present paper, the cases $r=4$ and $r=5$ are handled by a hybrid method:
the explicit $r$-adic recurrence pattern is first identified mathematically,
and the resulting universal identities are then certified by the Walnut theorem
prover (see Appendix~\ref{app:dfao}). The recurrence structure is first identified,
then the full verification is automatic once this structure is in place.

\begin{theorem}\label{thm:r4}
There are four sequences $\varepsilon_0, \varepsilon_1, \varepsilon_2, \varepsilon_3$ such that for all $n\ge 5$, $\varepsilon_i(n)\in\{0,1\}$, and
  \begin{align}
    a(4n)   &= a(n-3)+a(n-2)+a(n-1)+a(n)   + 1 + \varepsilon_0(n), \label{eq:r4a}\\
    a(4n+1) &= a(n-2)+a(n-1)+a(n)+a(n+1)   + \varepsilon_1(n),     \label{eq:r4b}\\
    a(4n+2) &= a(n-2)+a(n-1)+a(n)+a(n+1)   + \varepsilon_2(n),     \label{eq:r4c}\\
    a(4n+3) &= a(n-1)+a(n)+a(n+1)+a(n+2)   - 1 + \varepsilon_3(n). \label{eq:r4d}
  \end{align}
  The correction sequences $\varepsilon_i$ satisfy explicit $4$-adic recurrences
  (Theorem~\ref{thm:r4eps}); in particular, each $\varepsilon_i$ is $4$-automatic,
  and $(a(n))$ is $4$-regular.
\end{theorem}

\begin{remark}
  The windows slide one step to the right as the residue increases from $0$ to $3$:
  the window for residue $0$ is $(n-3,n-2,n-1,n)$, for residues $1$ and $2$ it is
  $(n-2,n-1,n,n+1)$, and for residue $3$ it is $(n-1,n,n+1,n+2)$.
  This is the same half-step sliding pattern as in the triadic case.
\end{remark}

\begin{proof}
  We define $\varepsilon_i(n)$ by the identities
  \eqref{eq:r4a}--\eqref{eq:r4d}.
  That $\varepsilon_i(n)\in\{0,1\}$ for all $n\ge 5$ and that each $\varepsilon_i$
  is $4$-automatic follows from Theorem~\ref{thm:r4eps} below, which
  establishes explicit $4$-adic recurrences for the quadruple
  $(\varepsilon_0,\varepsilon_1,\varepsilon_2,\varepsilon_3)$.
\end{proof}

\subsection{Automaticity of the correction sequences}

The four correction sequences satisfy explicit $4$-adic recurrences, which we establish next.

\begin{theorem}\label{thm:r4eps}
The correction sequences $\varepsilon_0,\varepsilon_1,\varepsilon_2,\varepsilon_3$
satisfy the following initial values for $5\le n\le 23$:
\[
\begin{array}{c|ccccccccccccccccccc}
n & 5&6&7&8&9&10&11&12&13&14&15&16&17&18&19&20&21&22&23\\ \hline
\varepsilon_0(n)&0&1&1&1&1&1&1&1&1&1&1&1&0&1&0&1&1&0&1\\
\varepsilon_1(n)&1&0&0&0&0&0&0&0&0&0&0&0&1&0&1&0&0&1&0\\
\varepsilon_2(n)&1&1&1&1&1&1&1&1&1&1&1&0&1&0&1&1&0&1&0\\
\varepsilon_3(n)&0&0&0&0&0&0&0&0&0&0&0&1&0&1&0&0&1&0&1
\end{array}
\]
Moreover, for every integer $m\ge 6$ the following $4$-adic recurrences hold:
\begin{align}
\varepsilon_0(4m)   &= (1-\varepsilon_0(m))\,\varepsilon_0(m+1)+\varepsilon_0(m)\,\varepsilon_0(m-1), \label{eq:e0_0}\\
\varepsilon_0(4m+1) &= (1-\varepsilon_0(m))\,\varepsilon_0(m-1)+\varepsilon_0(m)\,\varepsilon_0(m+1), \label{eq:e0_1}\\
\varepsilon_0(4m+2) &= \varepsilon_0(m), \label{eq:e0_2}\\
\varepsilon_0(4m+3) &= \varepsilon_0(m+1), \label{eq:e0_3}
\end{align}
\begin{align}
\varepsilon_1(4m)   &= \varepsilon_2(m)\,\bigl(\varepsilon_1(m-1)\oplus \varepsilon_1(m+1)\bigr), \label{eq:e1_0}\\
\varepsilon_1(4m+1) &= 1-\varepsilon_2(m), \label{eq:e1_1}\\
\varepsilon_1(4m+2) &= \varepsilon_1(m), \label{eq:e1_2}\\
\varepsilon_1(4m+3) &= \varepsilon_1(m+1), \label{eq:e1_3}
\end{align}
\begin{align}
\varepsilon_2(4m)   &= (1-\varepsilon_0(m-1))\,\varepsilon_1(m-1)+\varepsilon_0(m-1)\,(1-\varepsilon_1(m-1))\,\varepsilon_2(m), \label{eq:e2_0}\\
\varepsilon_2(4m+1) &= \varepsilon_2(m-1), \label{eq:e2_1}\\
\varepsilon_2(4m+2) &= \varepsilon_2(m), \label{eq:e2_2}\\
\varepsilon_2(4m+3) &= 1-\varepsilon_1(m), \label{eq:e2_3}
\end{align}
\begin{align}
\varepsilon_3(4m)   &= \varepsilon_2(m-1)\,\bigl(\varepsilon_2(m)\equiv \varepsilon_1(m)\bigr), \label{eq:e3_0}\\
\varepsilon_3(4m+1) &= \varepsilon_3(m-1), \label{eq:e3_1}\\
\varepsilon_3(4m+2) &= \varepsilon_3(m), \label{eq:e3_2}\\
\varepsilon_3(4m+3) &= (1-\varepsilon_1(m+1))\,\bigl(\varepsilon_2(m)\equiv \varepsilon_1(m)\bigr), \label{eq:e3_3}
\end{align}
where $x\oplus y := x+y-2xy$ and $x\equiv y := 1-x-y+2xy$ denote XOR and XNOR
for bits $x,y\in\{0,1\}$.
In particular, $\varepsilon_i(n)\in\{0,1\}$ for all $n\ge 5$, and each $\varepsilon_i$
is $4$-automatic.
\end{theorem}

\begin{proof}
The initial values follow by direct evaluation of the definitions of
$\varepsilon_0,\varepsilon_1,\varepsilon_2,\varepsilon_3$ in Theorem~\ref{thm:r4}
on the explicit initial segment of $a$.

The $4$-adic recurrences \eqref{eq:e0_0}--\eqref{eq:e3_3} were identified
by applying the defining equation \eqref{eq:r4def} at $n=4m,4m+1,4m+2,4m+3$
and expanding twice via \eqref{eq:r4a}--\eqref{eq:r4d}: all $a(\cdot)$-terms
cancel, leaving Boolean polynomial identities among the bits
$\varepsilon_j(m+\delta)$ with $|\delta|\le 1$.
Two illustrative derivations are given in Appendix~\ref{app:r4-deriv}
to display the elimination mechanism.
The complete system \eqref{eq:e0_0}--\eqref{eq:e3_3} is then certified
for all~$n$ by the Walnut theorem prover~\cite{Shallit2022}
(see Appendix~\ref{app:dfao} for the DFAO data);
Walnut also verifies the $\{0,1\}$-valuedness of all four sequences
and hence their $4$-automaticity.
\end{proof}

\begin{remark}
The initial values $\varepsilon_i(n)$ for $5\le n\le 23$ are obtained by
direct computation from the first $23$ terms of $a_4$. The Boolean
recurrences \eqref{eq:e0_0}--\eqref{eq:e3_3} then result from
eliminating the $a(\cdot)$-terms via the defining equation; two
representative cases are worked out in Appendix~\ref{app:r4-deriv}.
\end{remark}

\begin{corollary}\label{cor:r4regular}
The almost Golomb sequence of order $4$ is $4$-regular.
\end{corollary}

\begin{proof}
Theorem~\ref{thm:r4} expresses $a(4n+i)$ as a $\Z$-linear combination of
finitely many shifts of $a$ plus the $4$-automatic sequences $\varepsilon_i$
of Theorem~\ref{thm:r4eps}. By~\cite[Theorem~2.5]{AS92}, a sequence whose
$k$-kernel elements are all $\mathbb{Z}$-linear combinations of finitely many
fixed subsequences is $k$-regular; the recurrences in Theorem~\ref{thm:r4}
verify this condition for $k=4$, giving $(a(n))$ is $4$-regular.
\end{proof}

\section{The case $r=5$}\label{sec:r5}
%-------------------------------------------------------------------

For $r=5$, the proof follows the same hybrid pattern as for $r=4$:
one first identifies the correct quinary correction sequences and
residue-class formulas, then uses Walnut to certify the resulting
automatic identities, yielding a complete proof of the denesting
formulas and $5$-regularity.

\subsection{Two binary correction sequences}

For $n\ge 3$, define two correction sequences by
\begin{equation}\label{eq:r5_def}
  \varepsilon(n) := a(5n) - T_5(n) - 2,
  \qquad
  \eta(n) := T_5(n+2) - 1 - a(5n+2),
\end{equation}
where $T_5(m):=a(m)+a(m-1)+a(m-2)+a(m-3)+a(m-4)$ with $a(k)=0$ for $k<1$.

The two sequences $\varepsilon$ and $\eta$ are governed by the following $5$-adic recurrences.

\begin{theorem}\label{thm:r5eps}
The sequences $\varepsilon$ and $\eta$ take values in $\{0,1\}$,
satisfy $\varepsilon(n)\eta(n)=0$ for all $n\ge 3$, and are determined by
\[
  \varepsilon(1)=\varepsilon(2)=\varepsilon(3)=0,\quad \varepsilon(4)=1,
  \qquad
  \eta(1)=\eta(2)=0,\quad \eta(3)=1,\quad \eta(4)=0,
\]
together with the $5$-adic recurrences, valid for all $m\ge 4$:
\begin{align}
  \varepsilon(5m)   &= \varepsilon(m-1)\bigl(1-\varepsilon(m)\bigr), \label{eq:r5eps0}\\
  \varepsilon(5m+1) &= \varepsilon(m),\qquad \varepsilon(5m+3) = \varepsilon(m), \label{eq:r5eps13}\\
  \varepsilon(5m+2) &= \eta(m),\qquad\quad \varepsilon(5m+4) = \eta(m), \label{eq:r5eps24}\\
  \eta(5m)          &= \varepsilon(m),\qquad \eta(5m+2) = \varepsilon(m),\qquad
                       \eta(5m+4) = \varepsilon(m), \label{eq:r5eta024}\\
  \eta(5m+1)        &= \eta(m),\qquad\quad \eta(5m+3) = \eta(m). \label{eq:r5eta13}
\end{align}
In particular, the pair $U(n):=(\varepsilon(n),\eta(n))$ takes values in
$\{(0,0),(1,0),(0,1)\}$ and is $5$-automatic.
\end{theorem}

\begin{proof}
That $\varepsilon(n),\eta(n)\in\{0,1\}$ and $\varepsilon(n)\eta(n)=0$ holds for
$n\le 4$ by inspection.
The recurrences \eqref{eq:r5eps0}--\eqref{eq:r5eta13} were identified by applying
the defining equation $a(T_5(n))=n$ at $n=5m,\ldots,5m+4$, expanding each
$a(5m+j)$ in the window sum via \eqref{eq:r5_def}, and cancelling the $a(\cdot)$
terms. Two illustrative derivations are given in Appendix~\ref{app:r5-deriv}
to display the elimination mechanism.
The complete system of quinary recurrences
is then certified for all~$n$ by the Walnut theorem prover~\cite{Shallit2022}
(see Appendix~\ref{app:dfao}); Walnut also
verifies \eqref{eq:r5eps0}--\eqref{eq:r5eta13}, the $\{0,1\}$-valuedness
of $\varepsilon$ and $\eta$, and the disjointness identity
$\varepsilon(n)\eta(n)=0$ for all $n\ge 3$.

The right-hand side of each recurrence is either a constant, $\varepsilon(m)$,
$\eta(m)$, or $\varepsilon(m-1)(1-\varepsilon(m))$, hence lies in $\{0,1\}$.
Since $5m+d>m$ for $m\ge 1$, the recurrences define $(\varepsilon,\eta)$ on all
positive integers by induction. The disjointness $\varepsilon\eta=0$ is propagated
by induction: \eqref{eq:r5eps24} and \eqref{eq:r5eta024} exchange the roles of
$\varepsilon$ and $\eta$, while \eqref{eq:r5eps0} produces $1$ only when
$\varepsilon(m-1)=1$, hence when $\eta(m-1)=0$, so $\eta(5m)=\varepsilon(m)=0$.
Finally, $U(5m+d)$ is determined by $U(m)$ and $U(m-1)$, so the Allouche--Shallit
criterion~\cite{AS12} applies and $U$ is $5$-automatic.
\end{proof}

\begin{remark}
The five transitions of the pair $(\varepsilon(m),\eta(m))$ at consecutive
indices suffice to cover all cases in $a_5$; see Appendix~\ref{app:r5-tables}.
\end{remark}


\medskip
\begin{center}
\renewcommand{\arraystretch}{1.3}
\begin{tabular}{c|cccccccccc}
$n$ & 1 & 2 & 3 & 4 & 5 & 6 & 7 & 8 & 9 & 10 \\\hline
$\varepsilon(n)$ & 0 & 0 & 0 & 1 & 0 & 0 & 0 & 0 & 0 & 0 \\
$\eta(n)$ & 0 & 0 & 1 & 0 & 0 & 0 & 0 & 0 & 0 & 0 \\
\hline
$n$ & 11 & 12 & 13 & 14 & 15 & 16 & 17 & 18 & 19 & 20 \\\hline
$\varepsilon(n)$ & 0 & 0 & 1 & 0 & 1 & 0 & 1 & 0 & 1 & 0 \\
$\eta(n)$ & 0 & 1 & 0 & 1 & 0 & 1 & 0 & 1 & 0 & 1 \\
\end{tabular}
\end{center}
\medskip
\subsection{The denesting theorem}

Set $\sigma(n):=\varepsilon(n)+\eta(n)\in\{0,1\}$.
For residues $3$ and $4$ two further correctors appear, both expressible from
$(\varepsilon,\eta)$; their formulas are listed in Appendix~\ref{app:r5-tables}.

With these sequences in hand, the five denesting formulas are immediate.

\begin{theorem}\label{thm:r5}
For all $n\ge 3$, the almost Golomb sequence of order $5$ satisfies:
\begin{align}
  a(5n)   &= T_5(n)   + 2 + \varepsilon(n),   \label{eq:r5a0}\\
  a(5n+1) &= T_5(n+1) + 1 - \varepsilon(n),   \label{eq:r5a1}\\
  a(5n+2) &= T_5(n+2) - 1 - \eta(n),          \label{eq:r5a2}\\
  a(5n+3) &= T_5(n+3) - 4 + \theta(n),        \label{eq:r5a3}\\
  a(5n+4) &= T_5(n+4) + 2 + \varepsilon_4(n), \label{eq:r5a4}
\end{align}
where $\theta(n)\in\{0,1\}$ and $\varepsilon_4(n)\in\{-4,-3,-2\}$ are both
$5$-automatic, expressed as explicit functions of $U$ in
Appendix~\ref{app:r5-tables}. Consequently, $(a(n))$ is $5$-regular
by~\cite[Theorem~2.5]{AS92}.
\end{theorem}

\begin{proof}
Formulas \eqref{eq:r5a0} and \eqref{eq:r5a2} are the definitions
\eqref{eq:r5_def}. Formula \eqref{eq:r5a1} is derived in
Appendix~\ref{app:r5-deriv} by the same elimination used in Theorem~\ref{thm:r5eps}.
Formulas \eqref{eq:r5a3} and \eqref{eq:r5a4} are the definitions of $\theta$ and
$\varepsilon_4$.

The $5$-automaticity of $\theta$ and $\varepsilon_4$ follows from the explicit
tables in Appendix~\ref{app:r5-tables}: both are obtained from the $5$-automatic
sequence $U$ by a finite-state post-processing (the table defines a finite
transducer whose input is the pair $(U(m-1),U(m))$ or $(U(m),U(m+1))$).

The $5$-regularity of $(a(n))$ follows because \eqref{eq:r5a0}--\eqref{eq:r5a4}
express each subsequence $(a(5n+i))_{n\ge 0}$ as a $\Z$-linear combination of
finitely many shifts of $(a(n))$ plus $5$-automatic sequences; by
\cite[Theorem~2.5]{AS92} this implies $(a(n))$ is $5$-regular.
\end{proof}

%-------------------------------------------------------------------
\section{Non-convergence of $a(n)/n$}\label{sec:nonconv}
%-------------------------------------------------------------------

A common feature of all almost Golomb sequences is that $a(n)/n$ oscillates
between two explicit rational limits, failing to converge. This contrasts
sharply with Golomb's sequence, where $a(n)/n^{\varphi-1}\to\varphi^{2-\varphi}$
converges to a smooth constant.

\subsection{The case $r=2$}

\begin{theorem}\label{thm:ratio2}
  Let $(a(n))$ be the almost Golomb sequence of order $2$.
  \begin{enumerate}[label=(\roman*)]
    \item $a(n)/n$ does not converge: $a(2^k)/2^k = 3/4$ and
          $a(3\cdot 2^{k-1})/(3\cdot 2^{k-1}) = 2/3$ for all $k\ge 2$.
    \item The Cesàro means $C_N = \frac{1}{N}\sum_{n=1}^N \frac{a(n)}{n}$
          do not converge. They have two limit points:
          \[
            L_1 = \lim_{k\to\infty} C_{2^k}
                = \tfrac{3}{4} + \log\!\left(\tfrac{6^{3/4}}{4}\right),
            \qquad
            L_2 = \lim_{k\to\infty} C_{3\cdot 2^{k-1}}
                = \tfrac{2}{3} + \log\!\left(\tfrac{3^{2/3}}{2}\right).
          \]
  \end{enumerate}
\end{theorem}

\begin{proof}
  Part (i) follows immediately from Corollary~\ref{cor:pivots}.

  For part (ii), introduce the dyadic block average
  $A_k = 2^{-k}\sum_{n=2^k}^{2^{k+1}-1} a(n)/n$.
  Writing $n=2^k+j$ and $t=j/2^k\in[0,1)$, Proposition~\ref{prop:dyadic} gives
  \[
    \frac{a(n)}{n} =
    \begin{cases}
      \dfrac{\tfrac{3}{4}+\tfrac{t}{2}}{1+t} + O(2^{-k}) & 0\le t\le \tfrac{1}{2},\\[8pt]
      \dfrac{\tfrac{1}{2}+t}{1+t}                         & \tfrac{1}{2}<t<1.
    \end{cases}
  \]
  Summing as a Riemann integral gives $A_k\to\alpha$ where
  \[
    \alpha = \int_0^{1/2}\frac{\tfrac{3}{4}+\tfrac{t}{2}}{1+t}\,dt
           + \int_{1/2}^1  \frac{\tfrac{1}{2}+t}{1+t}\,dt
           = \tfrac{3}{4}+\log\!\left(\tfrac{6^{3/4}}{4}\right).
  \]
  Decomposing $C_{2^k}$ into dyadic blocks and using the fact that
  the weighted average $\sum_{j=0}^{k-1} 2^{j-k} A_j$ converges to $\alpha$
  whenever each $A_j\to\alpha$ (a standard Cesàro-type argument,
  since the weights $2^{j-k}$ are nonnegative and sum to $1-2^{-k}\to 1$)
  gives $\lim C_{2^k}=\alpha=L_1$.

  For $N_k=3\cdot 2^{k-1}$, write
  $C_{N_k} = \frac{2}{3}C_{2^k} + \frac{1}{3}D_k$
  where $D_k = 2^{1-k}\sum_{n=2^k+1}^{3\cdot 2^{k-1}} a(n)/n$.
  Since this sum lies in the first half of the dyadic block ($t\in[0,1/2]$),
  $D_k\to\beta$ with
  \[
    \beta = 2\int_0^{1/2}\frac{\tfrac{3}{4}+\tfrac{t}{2}}{1+t}\,dt
          = \tfrac{1}{2}+\tfrac{1}{2}\log\!\left(\tfrac{3}{2}\right).
  \]
  Hence $\lim C_{N_k} = \frac{2}{3}\alpha+\frac{1}{3}\beta = L_2$.
  Finally $L_1-L_2 = (1-\log(8/3))/12 > 0$ since $e>8/3$.
\end{proof}

The run structure of the sequence is also captured automatically.

\begin{proposition}
  The sequence $\Delta a(n) = a(n+1)-a(n) \in \{0,1\}$ is $2$-automatic.
\end{proposition}

\begin{proof}
  Proposition~\ref{prop:dyadic} shows that in the block $[2^k, 2^{k+1})$
  written as $n=2^k+j$:
  $\Delta a(n)=0$ when $j$ is odd and $j\le 2^{k-1}-1$
  (the ceiling-based formula produces repeated values at consecutive even/odd pairs);
  $\Delta a(n)=1$ for all $n\ge 3\cdot 2^{k-1}$ (the second half of the block
  has slope $1$).
  This is a regular language on the binary representation of $n$, so
  $\Delta a$ is $2$-automatic.
\end{proof}

The multiplicity sequence of values in $(a(n))$ is also $2$-automatic.

\begin{proposition}\label{prop:N2_multiplicity}
  Let $N_2(n)$ denote the number of times $n$ appears in the almost Golomb
  sequence of order $2$. Then $N_2(n)\in\{1,2\}$ for all $n\ge 1$, and the
  sequence $(N_2(n))_{n\ge 1}$ is $2$-automatic, satisfying the recurrences
  \begin{equation}\label{eq:N2rec}
    N_2(2n) = N_2(n),\qquad N_2(2n+1) = N_2(n+1)\qquad (n\ge 2),
  \end{equation}
  with initial values $N_2(1)=1$, $N_2(2)=2$, $N_2(3)=1$.
  Equivalently, $N_2$ is given by the single recursion
  \begin{equation}\label{eq:N2formula}
    N_2(n) = \begin{cases}
      2 - (n \bmod 2) & \text{if } n \le 3, \\
      N_2\!\bigl(\lceil n/2 \rceil\bigr) & \text{if } n \ge 4.
    \end{cases}
  \end{equation}
\end{proposition}

\begin{proof}
  From Proposition~\ref{prop:dyadic}, in the dyadic block $[2^k, 2^{k+1})$
  written as $n=2^k+j$ with $0\le j<2^k$:
  \begin{itemize}
    \item First half ($0\le j\le 2^{k-1}$): $a(n)=3\cdot 2^{k-2}+\lceil j/2\rceil$.
    The value $3\cdot 2^{k-2}$ appears once ($j=0$); each value in
    $[3\cdot 2^{k-2}+1,\,2^k]$ appears exactly twice (two consecutive values of $j$).
    \item Second half ($2^{k-1}<j<2^k$): $a(n)=n-2^{k-1}$, so
    each value in $[2^k+1,\,3\cdot 2^{k-1}-1]$ appears exactly once.
  \end{itemize}
  Hence $N_2(v)=2$ if and only if $v\in(3\cdot 2^{k-2},\,2^k]$ for some $k\ge 1$.

  We check that \eqref{eq:N2formula} agrees with this. The map
  $v\mapsto\lceil v/2\rceil$ descends until reaching $\{1,2,3\}$.
  Indeed, $v\in(3\cdot 2^{k-2},\,2^k]$ iff the descent
  $v\to\lceil v/2\rceil\to\lceil\lceil v/2\rceil/2\rceil\to\cdots$
  terminates at $2$: if $v$ is in the upper quarter of its dyadic block,
  $\lceil v/2\rceil$ lands in the upper quarter of the next smaller block,
  and the property is preserved until reaching $2\in(3/2,2]$.
  Conversely if $v$ is in the lower three-quarters, the descent
  eventually reaches $1$ or $3$, both of which give $N_2=1$ via~\eqref{eq:N2formula}.

  The recurrences~\eqref{eq:N2rec} are a direct consequence of
  \eqref{eq:N2formula}: for $n\ge 4$,
  $N_2(2n)=N_2(\lceil 2n/2\rceil)=N_2(n)$, and
  $N_2(2n+1)=N_2(\lceil(2n+1)/2\rceil)=N_2(n+1)$.
  Since $N_2$ takes values in $\{1,2\}$, it is bounded.
  The recurrences show that every subsequence in the $2$-kernel of $N_2$
  is a $\mathbb{Z}$-linear combination of
  $(N_2(n))_{n\ge 1}$, $(N_2(n+1))_{n\ge 1}$, and the constant sequence~$1$.
  Hence $N_2$ is $2$-regular. Since $N_2$ takes values in the finite
  set $\{1,2\}$, $2$-regularity implies $2$-automaticity~\cite{AS92}.
\end{proof}

%-------------------------------------------------------------------
\subsection{The case $r=3$}

To identify the two pivotal subsequences, we first locate them outside the correction set.

\begin{lemma}\label{lem:gaps}
  For $m\in\{5,8\}$ and all $k\ge 1$:\; $\varepsilon(m\cdot 3^k)=\varepsilon(m\cdot 3^k-1)=0$.
\end{lemma}

\begin{proof}
  Recall that $I_0=\{5\}$ and $I_k=[\ell_k,u_k]\cap\mathbb{Z}$ with
  $\ell_k=\lceil(11\cdot 3^k-1)/2\rceil$ and $u_k=\lfloor(13\cdot 3^k-3)/2\rfloor$
  for $k\ge 1$.
  The intervals are ordered by scale: $u_{k-1}<\ell_k$ for all $k\ge 1$, since
  $u_{k-1}=(13\cdot 3^{k-1}-3)/2$ and $\ell_k=(11\cdot 3^k-1)/2$ give
  $\ell_k - u_{k-1} = (11\cdot 3^k - 13\cdot 3^{k-1} + 2)/2 = (20\cdot 3^{k-1}+2)/2 > 0$.
  Hence the $I_k$ are pairwise disjoint, and $n\in\mathcal{I}$
  iff $n\in I_k$ for a unique $k$.

  For $m=5$ and $k\ge 1$, we show $5\cdot 3^k$ lies strictly between $I_{k-1}$ and $I_k$.
  On the left: $5\cdot 3^k > u_{k-1} = (13\cdot 3^{k-1}-3)/2$
  since $10\cdot 3^k = 30\cdot 3^{k-1} > 13\cdot 3^{k-1} - 3$.
  On the right: $5\cdot 3^k < \ell_k = (11\cdot 3^k-1)/2$
  since $10\cdot 3^k < 11\cdot 3^k - 1$, i.e.\ $1 < 3^k$, which holds for $k\ge 1$.
  Hence $5\cdot 3^k\notin\mathcal{I}$, and since $5\cdot 3^k - 1 < 5\cdot 3^k < \ell_k$
  and $5\cdot 3^k - 1 > u_{k-1}$, we get $5\cdot 3^k - 1\notin\mathcal{I}$ as well.

  For $m=8$ and $k\ge 1$, we show $8\cdot 3^k$ lies strictly between $I_k$ and $I_{k+1}$.
  On the left: $8\cdot 3^k > u_k = (13\cdot 3^k-3)/2$
  since $16\cdot 3^k > 13\cdot 3^k - 3$.
  On the right: $8\cdot 3^k < \ell_{k+1} = (11\cdot 3^{k+1}-1)/2$
  since $16\cdot 3^k < 33\cdot 3^k - 1$.
  Hence $8\cdot 3^k\notin\mathcal{I}$, and since $u_k < 8\cdot 3^k - 1 < 8\cdot 3^k < \ell_{k+1}$,
  we get $8\cdot 3^k - 1\notin\mathcal{I}$.
  Therefore $\varepsilon(5\cdot 3^k)=\varepsilon(5\cdot 3^k-1)
  =\varepsilon(8\cdot 3^k)=\varepsilon(8\cdot 3^k-1)=0$.
\end{proof}

\begin{theorem}\label{thm:ratio3}
  Let $(a(n))$ be the almost Golomb sequence of order $3$.
  For all $k\ge 1$:
  \[
    a(5\cdot 3^k) = 8\cdot 3^{k-1}, \quad
    \frac{a(5\cdot 3^k)}{5\cdot 3^k} = \frac{8}{15};
    \qquad
    a(8\cdot 3^k) = 15\cdot 3^{k-1}, \quad
    \frac{a(8\cdot 3^k)}{8\cdot 3^k} = \frac{5}{8}.
  \]
  Thus $a(n)/n$ does not converge.
\end{theorem}

\begin{proof}
  We propagate the local pattern around $5\cdot 3^k$ by induction on $k$.
  The base case $k=1$ is verified directly (pattern around $15$).
  Assuming for some $k\ge 1$:
  \[
    a(5\cdot 3^k-2)=8\cdot 3^{k-1}-1,\;\;
    a(5\cdot 3^k-1)=8\cdot 3^{k-1},\;\;
    a(5\cdot 3^k)  =8\cdot 3^{k-1},\;\;
    a(5\cdot 3^k+1)=8\cdot 3^{k-1}+1,
  \]
  set $n=5\cdot 3^k$. By Lemma~\ref{lem:gaps}, $\varepsilon(n)=\varepsilon(n-1)=0$.
  Applying \eqref{eq:r3a}--\eqref{eq:r3c} (and the analogues with $n-1$):
  \begin{align*}
    a(3n-2) &= a(n-2)+a(n-1)+a(n)   = 8\cdot 3^k-1,\\
    a(3n-1) &= a(n-1)+a(n)+a(n+1)-1 = 8\cdot 3^k,\\
    a(3n)   &= a(n-2)+a(n-1)+a(n)+1 = 8\cdot 3^k,\\
    a(3n+1) &= a(n-1)+a(n)+a(n+1)   = 8\cdot 3^k+1,
  \end{align*}
  which is the same pattern at rank $k+1$ around $3n=5\cdot 3^{k+1}$.

  For the family $a(8\cdot 3^k)=15\cdot 3^{k-1}$, we propagate the local pattern
  \[
    a(8\cdot 3^k-2)=15\cdot 3^{k-1}-1,\quad
    a(8\cdot 3^k-1)=a(8\cdot 3^k)=15\cdot 3^{k-1},\quad
    a(8\cdot 3^k+1)=15\cdot 3^{k-1}+1.
  \]
  The base case $k=1$ is read off directly: $a(22)=14$, $a(23)=a(24)=15$, $a(25)=16$.
  Assuming the pattern at rank $k$, set $n=8\cdot 3^k$.
  By Lemma~\ref{lem:gaps}, $\varepsilon(n)=\varepsilon(n-1)=0$.
  Applying \eqref{eq:r3a}--\eqref{eq:r3c} at $n-1$ and $n$:
  \begin{align*}
    a(3n-2) &= a(n-2)+a(n-1)+a(n)   = 15\cdot 3^k-1,\\
    a(3n-1) &= a(n-1)+a(n)+a(n+1)-1 = 15\cdot 3^k,\\
    a(3n)   &= a(n-2)+a(n-1)+a(n)+1 = 15\cdot 3^k,\\
    a(3n+1) &= a(n-1)+a(n)+a(n+1)   = 15\cdot 3^k+1,
  \end{align*}
  which is the same pattern at rank $k+1$ around $3n=8\cdot 3^{k+1}$.
  Since $8/15\ne 5/8$, the ratio $a(n)/n$ does not converge.
\end{proof}

%-------------------------------------------------------------------
\subsection{The case $r=4$}

The quaternary case exhibits the same oscillatory behavior as the triadic case,
but the two pivotal subsequences are now identified via the affine recurrence
$u_{k+1}=4u_k - 2$.

\begin{theorem}\label{thm:ratio4}
  Let $(a(n))$ be the almost Golomb sequence of order $4$.
  The sequence $a(n)/n$ does not converge. Setting $A_k=a(4^k)$,
  the recurrence
  \begin{equation}\label{eq:r4-rec}
    A_{k+1} = 4A_k - 2
  \end{equation}
  holds for all $k\ge 3$, with general solution $A_k = C\cdot 4^k + \tfrac{2}{3}$.
  The two families
  \[
    \frac{a(4^k)}{4^k} \to \frac{25}{48}, \qquad
    \frac{a(7\cdot 4^k)}{7\cdot 4^k} \to \frac{10}{21}
  \]
  have distinct limits, so $a(n)/n$ has at least two limit points.
\end{theorem}

\begin{proof}
  Since $4^k=[1\,0^k]_4$ and $7\cdot 4^k=[13\,0^k]_4$ in base~$4$,
  Lemma~\ref{lem:auto-geometric} implies that the correction patterns along
  these two families are ultimately periodic.
  Machine evaluation of the certified recurrence system
  of Theorem~\ref{thm:r4eps} yields, from $k=3$ onward,
  \[
    (\varepsilon_0(4^k),\,\varepsilon_1(4^k),\,\varepsilon_2(4^k),\,\varepsilon_3(4^k))
    =(1,0,0,1),
  \]
  together with the fixed local neighbourhood pattern
  \[
    a(4^k+j) = A_k + \delta(j) \qquad (-4\le j\le 4),
  \]
  where
  \[
    \delta(j) = \begin{cases}
      -2 & j\in\{-4,-3\},\\
      -1 & j\in\{-2,-1\},\\
      \phantom{-}0 & j\in\{0,1,2\},\\
      +1 & j\in\{3,4\}.
    \end{cases}
  \]
  Substituting into \eqref{eq:r4a} with $n=4^k$ gives
  $A_{k+1} = 4A_k - 2$.
  With $A_3 = a(64) = 34$, the solution is
  $A_k = \tfrac{25}{48}\cdot 4^k + \tfrac{2}{3}$.
  The same evaluation along $7\cdot 4^k$, with $B_k:=a(7\cdot 4^k)$, gives
  $B_k/(7\cdot 4^k)\to 10/21$.
  Since $25/48 \ne 10/21$, the ratio $a(n)/n$ does not converge.
  (The local neighbourhood patterns can be verified by
  iterating the certified recurrences of Theorem~\ref{thm:r4eps}.)
\end{proof}

\begin{remark}
  The affine recurrence $u_{k+1}=4u_k-2$ (with fixed point $2/3$) contrasts with
  the purely multiplicative recurrences of orders $2$ and $3$.
  Consequently the two limiting ratios $25/48$ and $10/21$ do not have the
  ``$p/q$ and $q/p$'' symmetry seen in the triadic case ($8/15$ and $5/8$).
\end{remark}

%-------------------------------------------------------------------
\subsection{The case $r=5$}

\begin{theorem}\label{thm:ratio5}
The ratio $a(n)/n$ for the almost Golomb sequence of order $5$ does not converge.
Setting $A_k=a(5^k)$ and $B_k=a(2\cdot 5^k)$, one has for all $k\ge 2$:
\begin{equation}\label{eq:r5rec}
  A_{k+1} =
  \begin{cases}
    5A_k - 1, & k\ \text{even},\\
    5A_k - 4, & k\ \text{odd},
  \end{cases}
  \qquad
  B_{k+1} = 5B_k - 1.
\end{equation}
Consequently
\[
  \lim_{\substack{k\to\infty\\k\text{ even}}} \frac{a(5^k)}{5^k} = \frac{93}{200},
  \qquad
  \lim_{k\to\infty}\frac{a(2\cdot 5^k)}{2\cdot 5^k} = \frac{87}{200},
\]
and these two limits are distinct.
\end{theorem}

\begin{proof}
Since $5^k=[1\,0^k]_5$ and $2\cdot 5^k=[2\,0^k]_5$ in base~$5$,
Lemma~\ref{lem:auto-geometric} implies that the correction patterns along
these families are ultimately periodic.
Machine evaluation of the certified recurrence system of
Theorem~\ref{thm:r5eps} yields, for $k\ge 2$, the alternating pattern
$U(5^k)\in\{(1,0),(0,1)\}$ with period~$2$, and a constant correction pattern
along $2\cdot 5^k$.
The same evaluation gives the corresponding local neighbourhood patterns of $a$,
from which the affine recurrences~\eqref{eq:r5rec} follow by substitution into
\eqref{eq:r5a0}--\eqref{eq:r5a4}.
For even $k$, $A_k/5^k \to 93/200$.
For $B_k$, $B_k/(2\cdot 5^k)\to 87/200$.
Since $93/200\ne 87/200$, the ratio $a(n)/n$
has at least two limit points.
(The local patterns follows by iterating the certified recurrences of Theorem~\ref{thm:r5eps}.)
\end{proof}

%-------------------------------------------------------------------

\subsection{Summary}

\medskip
\begin{center}
\renewcommand{\arraystretch}{1.6}
\begin{tabular}{c|l|c|l|c}
$r$ & Pivotal family 1 & Limit & Pivotal family 2 & Limit \\\hline
$2$ & $2^k$ & $\dfrac{3}{4}$ & $3\cdot 2^{k-1}$ & $\dfrac{2}{3}$ \\[4pt]
$3$ & $5\cdot 3^k$ & $\dfrac{8}{15}$ & $8\cdot 3^k$ & $\dfrac{5}{8}$ \\[4pt]
$4$ & $4^k$ & $\dfrac{25}{48}$ & $7\cdot 4^k$ & $\dfrac{10}{21}$ \\[4pt]
$5$ & $5^k$\;($k$ even) & $\dfrac{93}{200}$ & $2\cdot 5^k$ & $\dfrac{87}{200}$ \\
\end{tabular}
\end{center}
\medskip

In each case both limits lie in $(0,1)$, confirming linear growth $a(n)\asymp n$
(meaning $C_1 n \le a(n) \le C_2 n$ for positive constants $C_1,C_2$)
with an oscillating prefactor. For $r=2$, the Ces\`aro means also fail to converge
(Theorem~\ref{thm:ratio2}); the behavior for $r\ge 3$ is open
(see Section~\ref{sec:open}).

%-------------------------------------------------------------------
\section{Combinatorial interpretations}\label{sec:combinatorial}
%-------------------------------------------------------------------

The denesting formulas of Sections~\ref{sec:r2} and~\ref{sec:r3}
encode algebraic identities. In this section we extract from them
three combinatorial properties that reveal the
structure of the difference sequence $d(n)=a(n+1)-a(n)\in\{0,1\}$
in the cases $r=2$ and $r=3$.

\subsection{A second-bit characterization of multiplicities ($r=2$)}

Proposition~\ref{prop:N2_multiplicity} shows that the multiplicity
$N_2(m)$ depends only on the iterated ceiling map
$m\mapsto\lceil m/2\rceil$. We now translate this into a direct
binary characterization.

\begin{proposition}\label{prop:second-bit}
For every $m\ge 4$, write $m-1=(1\,b_2\,b_3\cdots b_\ell)_2$ with
$b_2$ denoting the second most significant bit. Then
\[
  N_2(m)=\begin{cases}
    2 & \text{if } b_2=1,\\
    1 & \text{if } b_2=0.
  \end{cases}
\]
\end{proposition}

\begin{proof}
By Proposition~\ref{prop:N2_multiplicity},
$N_2(m)=N_2(\lceil m/2\rceil)$ for $m\ge 4$.
Set $k=m-1\ge 3$ and write $k=(1\,b_2\cdots b_\ell)_2$.
The operation $m\mapsto\lceil m/2\rceil$ translates to
$k\mapsto\lfloor k/2\rfloor$, which erases the least significant bit.
After $\ell-2$ iterations the two-digit prefix
$(1\,b_2)_2$ is reached, corresponding to
$m'=k'+1$ where $k'\in\{2,3\}$.
If $b_2=1$: the descent reaches $k'=3$, hence $m'=4$, $N_2(4)=N_2(2)=2$.
If $b_2=0$: the descent reaches $k'=2$, hence $m'=3$, $N_2(3)=1$.
\end{proof}

\begin{remark}
The almost Golomb sequence of order $2$ is thus the unique nondecreasing
sequence of positive integers in which each value $m\ge 4$ is repeated
twice if the second most significant bit of $(m-1)$ is $1$, and once otherwise.
The asymptotic density of values appearing twice is
$\lim_{N\to\infty}|\{m\le N:N_2(m)=2\}|/N=1/2$,
which governs the oscillation of $a(n)/n$ between $2/3$ and $3/4$.
\end{remark}

\subsection{A Boolean cellular automaton ($r=2$)}

The denesting theorem (Theorem~\ref{thm:r2}) translates into a
remarkably simple rule for the difference sequence.

\begin{proposition}\label{prop:AND-automaton}
The difference sequence $d(n)=a(n+1)-a(n)\in\{0,1\}$ of the almost
Golomb sequence of order $2$ satisfies, for all $n\ge 2$:
\begin{equation}\label{eq:AND}
  d(2n)=1,\qquad d(2n+1)=d(n)\cdot d(n+1).
\end{equation}
\end{proposition}

\begin{proof}
From Theorem~\ref{thm:r2} (valid for $n\ge 2$):
\begin{align*}
  d(2n) &= a(2n+1)-a(2n) = \bigl[a(n)+a(n+1)\bigr]-\bigl[a(n)+a(n+1)-1\bigr]=1,\\
  d(2n+1) &= a(2n+2)-a(2n+1)
           = \bigl[a(n+1)+a(n+2)-1\bigr]-\bigl[a(n)+a(n+1)\bigr]\\
           &= d(n+1)+d(n)-1.
\end{align*}
Since $d\in\{0,1\}$ and $N_2(m)\in\{1,2\}$
(Proposition~\ref{prop:N2_multiplicity}), run lengths are at most~$2$,
so $d(n)+d(n+1)\ge 1$ for all $n\ge 2$.
Hence $d(n)+d(n+1)-1\in\{0,1\}$, and the identity
$x+y-1=xy$ holds for all $x,y\in\{0,1\}$ with $x+y\ge 1$.
\end{proof}

\begin{remark}
Formula~\eqref{eq:AND} shows that $d$ is generated by a multiplicative
cellular automaton: every even position carries the value $1$,
and every odd position is the Boolean AND of its two ``parent'' values
at positions $n$ and $n+1$.
The sequence $a(n)=\sum_{k=0}^{n-1}d(k)$ is thus the
accumulation function of this automaton.
\end{remark}

\subsection{A palindromic substitution rule ($r=3$)}

For $r=3$, the triadic structure of the difference sequence
admits a palindromic decomposition.

\begin{proposition}\label{prop:palindromic}
For every $n\ge 3$, define the triadic block
$B_n=(d(3n-2),\,d(3n-1),\,d(3n))$.
Then:
\begin{enumerate}[label=(\roman*)]
  \item $B_n$ is palindromic: $d(3n-2)=d(3n)$.
  \item The sum $d(3n-2)+d(3n-1)+d(3n)=L_n$.
  \item The block $B_n$ is determined by $L_n$ as follows:
    \[
      B_n=\begin{cases}
        (0,1,0) & \text{if } L_n=1,\\
        (1,0,1) & \text{if } L_n=2,\\
        (1,1,1) & \text{if } L_n=3.
      \end{cases}
    \]
    In particular, $L_n=3$ if and only if $\varepsilon(n-1)=1$.
\end{enumerate}
\end{proposition}

\begin{proof}
Using the denesting formulas \eqref{eq:r3a}--\eqref{eq:r3c}
at ranks $n-1$ and $n$:
\begin{align*}
  d(3n-2) &= a(3(n{-}1){+}2)-a(3(n{-}1){+}1) \\
          &= \bigl[a(n{-}1)+a(n)+a(n{+}1)-1-\varepsilon(n{-}1)\bigr]
             - \bigl[a(n{-}2)+a(n{-}1)+a(n)\bigr]\\
          &= a(n{+}1)-a(n{-}2)-1-\varepsilon(n{-}1) = L_n-1-\varepsilon(n{-}1).
\end{align*}
By the same method:
\begin{align*}
  d(3n-1) &= a(3n)-a(3(n{-}1){+}2) \\
          &= \bigl[a(n{-}2)+a(n{-}1)+a(n)+1+\varepsilon(n{-}1)\bigr]\\
          &\quad - \bigl[a(n{-}1)+a(n)+a(n{+}1)-1-\varepsilon(n{-}1)\bigr]\\
          &= 2-L_n+2\varepsilon(n{-}1).
\end{align*}
And:
\[
  d(3n) = a(3n{+}1)-a(3n) = L_n-1-\varepsilon(n{-}1) = d(3n{-}2).
\]
This proves (i). The sum is
$2(L_n-1-\varepsilon(n{-}1))+(2-L_n+2\varepsilon(n{-}1))=L_n$,
proving (ii).

For (iii), the invariant~\eqref{eq:r3-Linv} gives $L_n\le 2+\varepsilon(n{-}1)$
with $L_n=3\iff\varepsilon(n{-}1)=1$, so:
if $L_n=1$ (hence $\varepsilon(n{-}1)=0$): $d(3n{-}2)=d(3n)=0$, $d(3n{-}1)=1$;
if $L_n=2$ (hence $\varepsilon(n{-}1)=0$): $d(3n{-}2)=d(3n)=1$, $d(3n{-}1)=0$;
if $L_n=3$ (hence $\varepsilon(n{-}1)=1$): $d(3n{-}2)=d(3n)=1$, $d(3n{-}1)=1$.
\end{proof}

\begin{remark}\label{rem:oscillator}
Proposition~\ref{prop:palindromic} explains the oscillatory behavior of
$a(n)/n$ for $r=3$ as a delayed-feedback mechanism.
The infinite binary word $W=d(1)d(2)d(3)\cdots$ is self-describing:
the distance between consecutive $1$s in $W$ (at positions marking the
end of each run) equals the number of $1$s in the
sliding window of size $r=3$ ending at the corresponding index.
When the local density of $1$s is high (many short runs),
the rule dictates long future runs, creating a sparse zone;
when the window reaches this sparse zone, the count of $1$s drops,
forcing short runs and recreating a dense zone.
This perpetual oscillation is the source of the non-convergence of $a(n)/n$.
\end{remark}

\begin{remark}\label{rem:laplacian}
The propagation rules \eqref{eq:Lprop1}--\eqref{eq:Lprop3}
for the run lengths satisfy the relation
\[
  L_{3n-1}+L_{3n}+L_{3n+1}
  = L_{n-1}+L_n+L_{n+1}
    - \bigl(\varepsilon(n{-}2)-2\varepsilon(n{-}1)+\varepsilon(n)\bigr).
\]
The expression $f(n{-}1)-2f(n)+f(n{+}1)$ is the discrete Laplacian
$\Delta^2 f(n)$. Hence the total run length in a triadic block
is conserved up to a correction by $\Delta^2\varepsilon(n{-}1)$:
the points where $\varepsilon=1$ act as localized sources that
redistribute mass between adjacent runs, analogous to a discrete
diffusion process.
\end{remark}

%-------------------------------------------------------------------
\section{A conjectural Golomb meta-structure}\label{sec:golomb_return}
%-------------------------------------------------------------------

The almost Golomb sequences of order $r$ were introduced as finite-memory variants
of Golomb's sequence, departing from it in growth rate, regularity, and oscillatory
behavior. Golomb's sequence appears to re-emerge within the almost Golomb
hierarchy through a boundary phenomenon that is strongly supported numerically
and can be reduced to two precise conjectures.

For an almost Golomb sequence $(a(n))$ of order $r$, let $N_r(n)$ denote
the number of times $n$ appears in the sequence.
In Golomb's sequence, the self-describing property gives
$N(n)=G(n)\sim C n^{\varphi-1}$, which grows without bound.
For finite $r$, by contrast, $N_r(n)$ is bounded since $a(n)\asymp n$.
For $r=2$ this is made precise in Proposition~\ref{prop:N2_multiplicity}:
$N_2(n)\in\{1,2\}$ and $(N_2(n))$ is $2$-automatic.

Numerical computation for $r=2,\ldots,10$ yields the following maximal
multiplicities $M(r)=\sup_n N_r(n)$ (verified for $n$ well inside the computed range):
\[
\begin{array}{c|ccccccccc}
  r & 2 & 3 & 4 & 5 & 6 & 7 & 8 & 9 & 10 \\\hline
  M(r) & 2 & 3 & 3 & 3 & 3 & 4 & 4 & 4 & 5
\end{array}
\]
Let $j_k$ denote the smallest $r$ such that $M(r) \ge k$.
Numerical computation up to $r=700$ yields
\[
  j_3=3,\; j_4=7,\; j_5=10,\; j_6=13,\; j_7=17,\; j_8=21,\;
  j_9=25,\; j_{10}=30,\;\ldots
\]
with differences $j_{k+1}-j_k = 4,3,3,4,4,4,5,5,5,6,6,6,6,7,\ldots$
(starting from $k=3$).
The first gap $j_4-j_3=4$ is anomalous; from $k=4$ onwards the
difference sequence coincides exactly with Golomb's sequence:
\[
  j_{k+1}-j_k = G(k), \qquad k\ge 4,
\]
where $(G(n))_{n\ge 1}$ is Golomb's sequence as defined in
Section~\ref{sec:golomb} (OEIS \oeis{A001462}).
This has been verified for all $64$ gaps up to $r=700$.

Equivalently, letting $S(k)=\sum_{n=1}^{k}G(n)$ denote the partial
sums of Golomb's sequence (OEIS \oeis{A001463}), one has the closed formula
\[
  j_k = S(k-1) + 2, \qquad k\ge 4,
\]
verified for all $29$ values up to $r=200$.
Recall that $S(k)$ is characterised by $G(S(k))=k$: it is the position
of the last occurrence of $k$ in Golomb's sequence (OEIS \oeis{A001463}).
The Prefix Conjecture (below) posits $a_r(r)=G(r-1)$ for $r\ge 3$,
which for $r\ge 4$ gives $N_r(r-1)=G(r-1)$ via the run identity;
deriving the full threshold law $j_k=S(k-1)+2$ then requires in
addition the Domination Lemma.

If Conjecture~\ref{conj:golomb_return} holds, then since
$S(k)\sim \varphi^{2-\varphi}k^\varphi/\varphi$, one obtains
\[
  M(r) \;\sim\; C\,r^{\varphi-1}, \qquad r\to\infty,
\]
where $\varphi-1\approx 0.618$ is the same exponent as in
Golomb's own asymptotic $a(n)\sim C n^{\varphi-1}$.
If the conjecture holds, we come full circle. The sequences defined by
truncating Golomb's global rule to a finite window of size $r$ would carry,
in their run-length structure, the fingerprint of the very sequence they
were derived from.

Numerical evidence is collected in Table~\ref{tab:golomb_meta}.

%-------------------------------------------------------------------
\begin{table}[ht]
\centering
\renewcommand{\arraystretch}{1.35}
\caption{Numerical evidence for the Golomb meta-structure conjecture.
\textbf{Top:} $M(r)=\sup_n N_r(n)$ for $r=2,\ldots,50$.
\textbf{Bottom:} threshold positions $j_k$ (smallest $r$ with $M(r)\ge k$),
their successive differences $j_{k+1}-j_k$, and the values $G(k)$ of Golomb's
sequence~\cite{Golomb66}.
The identity $j_{k+1}-j_k = G(k)$ holds for all $k\ge 4$;
it has been verified for $28$ consecutive values up to $r=200$,
and for $64$ values up to $r=700$.}
\label{tab:golomb_meta}
\smallskip

{\small
\begin{tabular}{c|ccccccccccccccccc}
$r$    &  2& 3& 4& 5& 6& 7& 8& 9&10&11&12&13&14&15&16&17&18 \\\hline
$M(r)$ &  2& 3& 3& 3& 3& 4& 4& 4& 5& 5& 5& 6& 6& 6& 6& 7& 7
\end{tabular}

\smallskip
\begin{tabular}{c|ccccccccccccccccc}
$r$    & 19&20&21&22&23&24&25&26&27&28&29&30&31&32&33&34&35 \\\hline
$M(r)$ &  7& 7& 8& 8& 8& 8& 9& 9& 9& 9& 9&10&10&10&10&10&11
\end{tabular}

\smallskip
\begin{tabular}{c|ccccccccccccccc}
$r$    & 36&37&38&39&40&41&42&43&44&45&46&47&48&49&50 \\\hline
$M(r)$ & 11&11&11&11&12&12&12&12&12&12&13&13&13&13&13
\end{tabular}
}

\bigskip

{\small
\begin{tabular}{c|c|cccccccccccccc}
$k$             &  3 &  4 &  5 &  6 &  7 &  8 &  9 & 10 & 11 & 12 & 13 & 14 & 15 & 16 \\\hline
$j_k$           &  3 &  7 & 10 & 13 & 17 & 21 & 25 & 30 & 35 & 40 & 46 & 52 & 58 & 64 \\\hline
$j_{k+1}{-}j_k$ &  4 &  3 &  3 &  4 &  4 &  4 &  5 &  5 &  5 &  6 &  6 &  6 &  6 &  7 \\\hline
$G(k)$          &  2 &  3 &  3 &  4 &  4 &  4 &  5 &  5 &  5 &  6 &  6 &  6 &  6 &  7 \\
\end{tabular}

\smallskip
\begin{tabular}{c|cccccccccccccc}
$k$             & 17 & 18 & 19 & 20 & 21 & 22 & 23 & 24 & 25 & 26 & 27 & 28 & 29 & 30 \\\hline
$j_k$           & 71 & 78 & 85 & 92 &100 &108 &116 &124 &133 &142 &151 &160 &169 &179 \\\hline
$j_{k+1}{-}j_k$ &  7 &  7 &  8 &  8 &  8 &  8 &  9 &  9 &  9 &  9 &  9 & 10 & 10 & 10 \\\hline
$G(k)$          &  7 &  7 &  7 &  8 &  8 &  8 &  8 &  9 &  9 &  9 &  9 &  9 & 10 & 10 \\
\end{tabular}
}
\end{table}
%-------------------------------------------------------------------

\begin{conjecture}\label{conj:golomb_return}
  For all $k\ge 4$, the threshold $j_k$ (smallest order $r$ at which
  $M(r)=\sup_n N_r(n)$ first reaches $k$) satisfies
  \[
    j_{k+1}-j_k = G(k) \qquad \text{and} \qquad j_k = S(k-1)+2,
  \]
  where $G$ is Golomb's sequence and $S(k)=\sum_{n=1}^{k}G(n)$.
\end{conjecture}

\subsection{A conditional reduction}

While a complete proof of Conjecture~\ref{conj:golomb_return} remains open,
the phenomenon can be reduced to two precise boundary conjectures.

Set $S_m^{(r)} = \sum_{i=0}^{r-1} a_r(m-i)$ (with $a_r(k)=0$ for $k\le 0$),
so that $a_r(S_m^{(r)})=m$ by the defining equation.
The following lemma provides the run-length identity used in the reduction.

\begin{lemma}\label{lem:run_identity}
  For all $r\ge 2$ and $m\ge 3$,
  \begin{equation}\label{eq:run_diff}
    N_r(m) = a_r(m+1) - a_r(m+1-r).
  \end{equation}
\end{lemma}

\begin{proof}
  By Theorem~\ref{thm:global_run}, the run of value $m$ is exactly
  $[S_m^{(r)},S_{m+1}^{(r)}-1]$, so $N_r(m)=S_{m+1}^{(r)}-S_m^{(r)}$.
  Telescoping: $S_{m+1}^{(r)}-S_m^{(r)}=\sum_{j=0}^{r-1}a_r(m+1-j)-\sum_{j=0}^{r-1}a_r(m-j)
  =a_r(m+1)-a_r(m+1-r)$.
\end{proof}

The meta-structure phenomenon reduces to the value of a single term:
for $r\ge 4$, Lemma~\ref{lem:run_identity} at $m=r-1$ gives
$N_r(r-1)=a_r(r)-a_r(0)=a_r(r)$.

\begin{conjecture}[Prefix Conjecture]\label{conj:prefix}
  For all $r\ge 3$,
  \[
    a_r(r) = G(r-1).
  \]
\end{conjecture}

The conjecture has been verified computationally for all $r\le 200$.
Note that for $n\le r$, the convention $a(k)=0$ for $k\le 0$ ensures that the
sliding window reduces to the full sum $\sum_{k=1}^{n}a(k)$, so the
first $r$ terms of $a_r$ depend only on the greedy rule with full memory.
The conjecture thus asserts that this common initial prefix, evaluated at $n=r$,
gives $G(r-1)$.

Combining Conjecture~\ref{conj:prefix} with Lemma~\ref{lem:run_identity}
(which requires $m\ge 3$, hence $r\ge 4$) yields the boundary identity
\begin{equation}\label{eq:threshold}
  N_r(r-1)=a_r(r)=G(r-1),\qquad r\ge 4.
\end{equation}
(For $r=3$, Conjecture~\ref{conj:prefix} gives $a_3(3)=G(2)=2$, but the run identity
does not apply at $m=2<3$; indeed $N_3(2)=3\ne 2$.)

\begin{conjecture}[Domination Lemma]\label{conj:domination}
  For all $r\ge 5$,
  \[
    M(r) = N_r(r-1).
  \]
\end{conjecture}

(For $r=4$ the identity fails: $M(4)=3$ but $N_4(3)=2$.)

A proof of the Domination Lemma (Conjecture~\ref{conj:domination}) would require showing that $N_r(m)\le G(r-1)$
for all $m\ge r$, or equivalently that $S_{m+1}^{(r)}-S_m^{(r)}\le G(r-1)$.
This is strongly supported by the following comparison of growth regimes:
in the stationary range ($m\gg r$), the run lengths are numerically seen
to grow slowly with $r$, whereas the boundary run length
$G(r-1)\sim C r^{\varphi-1}$ with $\varphi-1\approx 0.618$ grows strictly faster.
The initial peak at the boundary therefore permanently dominates the stationary regime.

Assuming both conjectures, the derivation of the threshold law is immediate:
$M(r)=G(r-1)$, so $j_k$ is the smallest $r$ with $G(r-1)\ge k$,
giving $j_k=S(k-1)+2$ and $j_{k+1}-j_k=G(k)$.

\section{Open questions and perspectives}\label{sec:open}
%-------------------------------------------------------------------

With the Golomb meta-structure reduced to two precise conjectures (the Prefix Conjecture and the Domination Lemma), several structural and analytic questions remain open. We conclude with six open problems.

\begin{enumerate}
  \item \textbf{The singular limit and phase transition ($r \to \infty$).}
        For any finite $r$, the sequence exhibits oscillatory linear growth
        $a(n) \asymp n$, whereas Golomb's sequence (the formal limit $r\to\infty$)
        grows sublinearly as $G(n) \sim C n^{\varphi-1}$.
        How does the sequence undergo this phase transition?
        Specifically, letting $L_{\max}(r) = \limsup_{n\to\infty} a(n)/n$, numerical
        evidence suggests that $L_{\max}(r)$ is strictly decreasing
        (values $3/4$, $5/8$, $25/48$, $93/200$ for $r=2,3,4,5$ respectively).
        Does $\lim_{r\to\infty} L_{\max}(r) = 0$?

  \item \textbf{Explicit denesting structure.}
        Theorem~\ref{thm:structural_regularity} proves $r$-regularity for all $r$.
        However, the DFAO produced is large. For each $r$, does the almost Golomb
        sequence admit an explicit denesting $a(rn+i)=T_r(n+s_i)+\delta_i+\varepsilon_i(n)$
        with $\varepsilon_i$ $r$-automatic and satisfying compact Boolean recurrences,
        as established for $r\le 5$? A proof of this finer structure would explain why
        Walnut produces automata with $O(r^2)$ states rather than the theoretical bound.

  \item \textbf{Ces\`aro means.}
        For $r=2$, Theorem~\ref{thm:ratio2} shows that the Ces\`aro means of
        $a(n)/n$ do not converge. For $r \ge 3$, do the block averages
        $\frac{1}{N}\sum_{n=1}^N a(n)/n$ converge, or do they exhibit
        fractal fluctuations periodic in $\log_r N$?

  \item \textbf{Automaticity of the full multiplicity sequence.}
        For $r=2$, $(N_2(n))$ is $2$-automatic
        (Proposition~\ref{prop:N2_multiplicity}). Does the $r$-regularity
        of $(a(n))$ guarantee that $(N_r(n))_{n\ge 1}$ is always $r$-automatic?

  \item \textbf{Combinatorial interpretations for $r\ge 4$.}
        Section~\ref{sec:combinatorial} provides a second-bit characterization
        and a Boolean cellular automaton for $r=2$, and a palindromic substitution
        rule for $r=3$. Do the cases $r\ge 4$ admit analogous combinatorial
        descriptions of the difference sequence, and can the delayed-feedback
        mechanism of Remark~\ref{rem:oscillator} be made quantitative?

  \item \textbf{The gap-$s$ variant.}
        Fix a gap $s\ge 1$ and consider the sequence defined by
        $a\bigl(a(n)+a(n-s)\bigr)=n$ for $n\ge s+1$, with $a(k)=0$ for $k\le 0$
        and $a(1)=1$.

        For $s=1$ this is the almost Golomb sequence of order $r=2$ studied in
        Section~\ref{sec:r2}.

        For $s=2$, the sequence begins
        $1,2,2,3,3,4,5,6,6,7,7,8,8,\ldots$
        and satisfies $a(n+1)-a(n)\in\{0,1\}$ for all $n\ge 1$.
        One can show that the sequence admits an exact divide-and-conquer
        structure modulo a binary correction sequence $\varepsilon(n)\in\{0,1\}$,
        whose support is
        \[
          \mathcal{I} = \{5,12\}\;\cup\;\bigcup_{k\ge 2}\mathcal{I}_k,
          \qquad
          \mathcal{I}_k = \bigl[5\cdot 2^k+2,\;6\cdot 2^k-2\bigr]\cap(4\Z+2),
        \]
        with $|\mathcal{I}_k|=2^{k-2}$, satisfying $\mathcal{I}_k=\{6\cdot 2^k-2\}\cup(2\cdot\mathcal{I}_{k-1}-2)$.
        Two families of exact values hold:
        \[
          a(7\cdot 2^k) = 5\cdot 2^k\ (k\ge 1),\quad
          a(5\cdot 2^k) = 7\cdot 2^{k-1}\ (k\ge 2),\quad
          a(3\cdot 2^k) = 17\cdot 2^{k-3}\ (k\ge 3),
        \]
        and $a(n)/n$ oscillates between the exact limits
        \[
          \liminf_{n\to\infty}\frac{a(n)}{n} = \frac{7}{10},\qquad
          \limsup_{n\to\infty}\frac{a(n)}{n} = \frac{5}{7}.
        \]
        This should be compared with the Stern diatomic sequence~\cite[Chap.\,16]{AS03}
        (\oeis{A002487}), which satisfies $s(2n)=s(n)$, $s(2n+1)=s(n)+s(n+1)$,
        is $2$-regular with differences in $\{0,1\}$, and admits no correction
        sequence because its denesting is exact.

        For $s=3$, the sequence $1,3,3,4,4,5,5,6,7,8,8,$\linebreak$9,10,11,11,\ldots$
        has differences in $\{0,1,2\}$ and numerical investigation reveals
        no stable divide-and-conquer recurrence at any base.
        This case appears to escape the $2$-regular framework entirely.

        It is an open problem to determine for which values of $s$ the
        gap-$s$ almost Golomb sequence is $2$-regular, and whether $s\in\{1,2\}$
        are the only cases admitting an exact correction-set description.
\end{enumerate}

%-------------------------------------------------------------------
\appendix
\section{Illustrative quaternary derivations}\label{app:r4-deriv}
%-------------------------------------------------------------------

This appendix presents two sample derivations from
Theorem~\ref{thm:r4eps}, to illustrate the elimination mechanism
that produces the Boolean recurrences.
All sixteen recurrences are obtained by the same procedure; the
complete system is certified by Walnut (Appendix~\ref{app:dfao}),
which constitutes the formal proof.

\subsection*{Notation}
For $m\ge 5$ set
\begin{align*}
  X_m &:= a(m-3)+a(m-2)+a(m-1)+a(m),\\
  Y_m &:= a(m-2)+a(m-1)+a(m)+a(m+1),\\
  Z_m &:= a(m-1)+a(m)+a(m+1)+a(m+2).
\end{align*}
The denesting formulas \eqref{eq:r4a}--\eqref{eq:r4d} then read
\begin{align*}
  a(4m)   &= X_m+1+\varepsilon_0(m), &\quad
  a(4m+1) &= Y_m+\varepsilon_1(m),\\
  a(4m+2) &= Y_m+\varepsilon_2(m), &\quad
  a(4m+3) &= Z_m-1+\varepsilon_3(m).
\end{align*}
Note the identities $Z_{m-1}=Y_m$ and $Y_{m-1}=X_m$ (both immediate from
the definitions by shifting $m$ by $-1$).

\subsection*{Derivation 1: \eqref{eq:e0_2}, easy case $\varepsilon_0(4m+2)=\varepsilon_0(m)$}

By definition of $\varepsilon_0$,
\[
  \varepsilon_0(4m+2) = a\bigl(4(4m+2)\bigr) - X_{4m+2} - 1,
\]
where $X_{4m+2} = a(4m-1)+a(4m)+a(4m+1)+a(4m+2)$.
Expanding each term via \eqref{eq:r4a}--\eqref{eq:r4d}:
\begin{align*}
  a(4m-1) &= a(4(m-1)+3) = Z_{m-1}-1+\varepsilon_3(m-1) = Y_m-1+\varepsilon_3(m-1),\\
  a(4m)   &= X_m+1+\varepsilon_0(m),\\
  a(4m+1) &= Y_m+\varepsilon_1(m),\\
  a(4m+2) &= Y_m+\varepsilon_2(m).
\end{align*}
Hence $X_{4m+2} = 3Y_m+X_m + (\varepsilon_3(m-1)+\varepsilon_0(m)+\varepsilon_1(m)+\varepsilon_2(m))$.

Now expand $a(4(4m+2))$ by noting $4(4m+2) = 4\cdot(4m+2)$.
Since $4m+2\equiv 2\pmod{4}$, equation \eqref{eq:r4c} gives
$a(4(4m+2)) = Y_{4m+2}+\varepsilon_2(4m+2)$.
The sum $Y_{4m+2} = a(4m)+a(4m+1)+a(4m+2)+a(4m+3)$ is given by
\begin{align*}
  Y_{4m+2} &= (X_m+1+\varepsilon_0(m))+(Y_m+\varepsilon_1(m))+(Y_m+\varepsilon_2(m))+(Z_m-1+\varepsilon_3(m))\\
            &= X_m+2Y_m+Z_m+(\varepsilon_0(m)+\varepsilon_1(m)+\varepsilon_2(m)+\varepsilon_3(m)).
\end{align*}
Substituting and using $Z_m = Y_m+(a(m+2)-a(m-2))$:
\begin{align*}
  \varepsilon_0(4m+2)
  &= Y_{4m+2}+\varepsilon_2(4m+2) - X_{4m+2} - 1\\
  &= \bigl[X_m+3Y_m+a(m+2)-a(m-2)\bigr]
     + \bigl[\varepsilon_0+\varepsilon_1+\varepsilon_2+\varepsilon_3\bigr](m)
     + \varepsilon_2(4m+2)\\
  &\quad - \bigl[3Y_m+X_m+\varepsilon_3(m-1)+\varepsilon_0(m)+\varepsilon_1(m)+\varepsilon_2(m)\bigr] - 1\\
  &= (a(m+2)-a(m-2)-1) + \varepsilon_3(m)-\varepsilon_3(m-1) + \varepsilon_2(4m+2).
\end{align*}
Since $\varepsilon_0(4m+2)\in\{0,1\}$ and $\varepsilon_2(4m+2)\in\{0,1\}$, and
$a(m+2)-a(m-2)\in\{2,3,4\}$ (unit increments, window of length $4$), direct
comparison of cases yields $a(m+2)-a(m-2)-1 = 1+\delta$ and
$\varepsilon_3(m)-\varepsilon_3(m-1)+\varepsilon_2(4m+2) = \varepsilon_0(m)-\delta$
for $\delta\in\{0,1\}$, collapsing to $\varepsilon_0(4m+2)=\varepsilon_0(m)$.
This identity is equivalent to \eqref{eq:e0_2} and is certified by Walnut for all~$n$.

\subsection*{Derivation 2: \eqref{eq:e1_1}, typical case $\varepsilon_1(4m+1)=1-\varepsilon_2(m)$}

The defining equation at $n=4m+1$ states $a(S_{4m+1})=4m+1$, where
\[
  S_{4m+1} = a(4m+1)+a(4m)+a(4m-1)+a(4m-2).
\]
Expanding:
\begin{align*}
  a(4m+1) &= Y_m+\varepsilon_1(m),\\
  a(4m)   &= X_m+1+\varepsilon_0(m),\\
  a(4m-1) &= Y_m-1+\varepsilon_3(m-1),\\
  a(4m-2) &= X_m+\varepsilon_2(m-1).
\end{align*}
Hence $S_{4m+1} = 2X_m+2Y_m + (\varepsilon_1(m)+\varepsilon_0(m)+\varepsilon_3(m-1)+\varepsilon_2(m-1))$.
Writing $E_m := \varepsilon_0(m)+\varepsilon_1(m)+\varepsilon_2(m-1)+\varepsilon_3(m-1)\in\{0,1,2,3,4\}$
and $F_m := a(m-2)+a(m-1)+a(m)$, one has $2X_m+2Y_m = 4F_m+2(a(m-3)+a(m+1))$,
so $S_{4m+1} \equiv 2(a(m-3)+a(m+1))+E_m \pmod{4}$.

The residue of $S_{4m+1}$ modulo $4$ determines which of \eqref{eq:r4a}--\eqref{eq:r4d}
applies at $S_{4m+1}$, and therefore which $\varepsilon_j$ appears when one writes
$a(S_{4m+1})=4m+1$. In each of the four sub-cases, imposing $a(S_{4m+1})=4m+1$
and using the known values of $\varepsilon_j(m+\delta)$ for $|\delta|\le 1$ yields
the single Boolean identity
\[
  \varepsilon_1(4m+1) = 1-\varepsilon_2(m).
\]
This is \eqref{eq:e1_1}. The identity is certified by Walnut for all~$n$.

\subsection*{How to obtain the remaining recurrences}
Each of the fourteen remaining identities in \eqref{eq:e0_0}--\eqref{eq:e3_3} is
obtained by the same two-step procedure: first expand $X_{4m+i}$ or $S_{4m+i}$ using
\eqref{eq:r4a}--\eqref{eq:r4d}, then expand $a(4(4m+i))$ or $a(S_{4m+i})$
according to the residue class. The result is always a Boolean polynomial
in $\varepsilon_j(m+\delta)$ with $|\delta|\le 1$, which simplifies to the
stated recurrence. All sixteen identities are certified by Walnut.

%-------------------------------------------------------------------
\section{Illustrative quinary derivations}\label{app:r5-deriv}
%-------------------------------------------------------------------

\subsection*{Notation}

For $m\ge 4$ set
\[
  W_m := a(m)+a(m-1)+a(m-2)+a(m-3)+a(m-4) = T_5(m).
\]
The denesting formulas \eqref{eq:r5a0}--\eqref{eq:r5a2} read
\[
  a(5m) = W_m+2+\varepsilon(m),\quad
  a(5m+1) = W_{m+1}+1-\varepsilon(m),\quad
  a(5m+2) = W_{m+2}-1-\eta(m).
\]

\subsection*{Derivation of \eqref{eq:r5a1}: $a(5m+1)=T_5(m+1)+1-\varepsilon(m)$}

The defining equation at $n=5m+1$ reads
$a(S)=5m+1$ where $S=a(5m+1)+a(5m)+a(5m-1)+a(5m-2)+a(5m-3)$.
Expand the four terms $a(5m),a(5m-1),a(5m-2),a(5m-3)$ using
\eqref{eq:r5a0}, \eqref{eq:r5a2} (at $n=m-1$), and the neighbouring residue
formulas. The deterministic part of $S$ equals $T_5(m+1)+(1-\varepsilon(m))+W_m$,
and after the second expansion via $a(S)=5m+1$ all $a(\cdot)$-terms cancel,
yielding \eqref{eq:r5a1}.

\subsection*{Derivation of \eqref{eq:r5eps0}: $\varepsilon(5m)=\varepsilon(m-1)(1-\varepsilon(m))$}

By definition $\varepsilon(5m)=a(25m)-T_5(5m)-2$.
Write $T_5(5m)=a(5m)+a(5m-1)+\cdots+a(5m-4)$ and expand each term
using the residue formulas at scale $m$:
\begin{align*}
  a(5m)   &= W_m+2+\varepsilon(m),&
  a(5m-1) &= W_m+1-\varepsilon(m-1)(1-\varepsilon(m-1)), \\
  a(5m-2) &= W_m-\eta(m-1)-1,&
  a(5m-3) &= W_m-2+\theta(m-1),\\
  a(5m-4) &= W_m+2+\varepsilon_4(m-1).
\end{align*}
Summing and subtracting $2$ produces $\varepsilon(5m)$ in terms of Boolean
expressions of $(\varepsilon(m-1),\eta(m-1))$; simplification using the
transition structure of $U$ yields
$\varepsilon(5m)=\varepsilon(m-1)(1-\varepsilon(m))$, i.e.~\eqref{eq:r5eps0}.

\subsection*{How to obtain the remaining recurrences in
Theorem~\ref{thm:r5eps}}

Each of the nine remaining identities
\eqref{eq:r5eps13}--\eqref{eq:r5eta13} is obtained by the same procedure:
apply the defining equation at $n=5m+d$ for the appropriate $d$,
expand the window sum via the five residue formulas, and cancel the
$a(\cdot)$-terms. The result is always a function of $(\varepsilon(m),\eta(m))$
or $(\varepsilon(m-1),\eta(m-1))$, which simplifies to the stated recurrence.

%-------------------------------------------------------------------
\section{Quinary correction tables}\label{app:r5-tables}
%-------------------------------------------------------------------

Write $U(n)=(\varepsilon(n),\eta(n))\in\{(0,0),(1,0),(0,1)\}$.
For $m\ge 3$ the only transitions $(U(m-1),U(m))$ that occur are:
\[
(0,0)\to(0,0),\quad (0,0)\to(0,1),\quad (0,1)\to(1,0),\quad
(1,0)\to(0,0),\quad (1,0)\to(0,1).
\]

\subsection*{Table for $\theta$}

The corrector $\theta(5m+d)\in\{0,1\}$ is determined by $(U(m-1),U(m))$ and
$d\in\{0,1,2,3,4\}$:
\[
\begin{array}{c|ccccc}
(U(m-1),U(m)) & d=0&1&2&3&4\\\hline
((0,0),(0,0)) & 1&1&1&1&1\\
((1,0),(0,0)) & 1&1&1&1&1\\
((0,0),(0,1)) & 0&1&0&1&0\\
((1,0),(0,1)) & 0&1&0&1&0\\
((0,1),(1,0)) & 1&0&1&0&1
\end{array}
\]

\subsection*{Table for $\varepsilon_4$}

For $d\in\{0,1,2,3\}$, set $\sigma(m)=\varepsilon(m)+\eta(m)$; then
\[
  \varepsilon_4(5m+d) =
  \begin{cases}
    -2-\sigma(m)-\varepsilon(m), & d\in\{0,2\},\\
    -2-\sigma(m)-\eta(m), & d\in\{1,3\}.
  \end{cases}
\]
For $d=4$, the value is determined by $(U(m),U(m+1))$:
\[
  \varepsilon_4(5m+4) =
  \begin{cases}
    -2,& (U(m),U(m+1))=((0,0),(0,0)),\\
    -3,& (U(m),U(m+1))\in\bigl\{((0,0),(0,1)),((0,1),(1,0)),((1,0),(0,0))\bigr\},\\
    -4,& (U(m),U(m+1))=((1,0),(0,1)).
  \end{cases}
\]
All entries are consistent with the $5$-automatic structure certified
by Walnut (Appendix~\ref{app:dfao}) and have been verified independently
for all $n\le 2000$.

%-------------------------------------------------------------------
\section{A menagerie of nested self-referential sequences}\label{app:nested}
%-------------------------------------------------------------------

The almost Golomb sequences studied in this paper belong to a broader landscape of
self-referential sequences admitting a denesting into base-$k$
divide-and-conquer formulas. We collect here the most relevant examples
for comparison, distinguishing two types.

\subsection*{Monotone sequences}

These are defined greedily: $a(n)$ is the \emph{smallest} integer $\ge a(n-1)$
consistent with the defining constraint. They are nondecreasing by construction,
unbounded, and hence $k$-regular (not $k$-automatic) for appropriate $k$;
see~\cite{AS92} and~\cite[Chapter~16]{AS03}.

\medskip
{\small
\begin{center}
\renewcommand{\arraystretch}{1.5}
\begin{tabular}{c|p{9cm}|c}
OEIS & \multicolumn{1}{c|}{Implicit rule \& denesting} & $k$-reg. \\\hline
\oeis{A003605} & $a(a(n))=3n$;\;
  $a(3n)=3a(n)$, $a(3n{+}1)=2a(n){+}a(n{+}1)$,
  $a(3n{+}2)=a(n){+}2a(n{+}1)$ & $3$ \\
\oeis{A079000} & $a(a(n))=2n+3$;\; $2$-adic formulas~\cite{CSV03,ARS05} & $2$ \\
\oeis{A080637} & $a(a(n))=2n+1$;\; $2$-adic formulas~\cite{CSV03,ARS05} & $2$ \\
\oeis{A394217} & $a(a(n){+}a(n{-}1))=n$;\;
  $a(2n)=a(n){+}a(n{+}1){-}1$,\; $a(2n{+}1)=a(n){+}a(n{+}1)$ & $2$ \\
\oeis{A394218} & $a(a(n){+}a(n{-}1){+}a(n{-}2))=n$;\;
  $3$-adic with correction $\one_\Ical$, \S\ref{sec:r3} & $3$ \\
\end{tabular}
\end{center}
}
\medskip

The last two rows are the almost Golomb sequences of orders $2$ and $3$ studied in
this paper, which generalize the $a(a(n))=kn$ family (studied systematically
in~\cite{ARS05}) by replacing the single
nested index $a(n)$ with a sliding window sum.
A related family of \emph{meta-automatic} sequences, combining
nested meta-Fibonacci recurrences with digit-based rules, is studied
in~\cite{CC25}.

\subsection*{Earliest (non-monotone) sequences}

These are defined differently: $a(n)$ is the \emph{smallest positive integer not
yet used} consistent with the constraint. They are \emph{not} nondecreasing in
general, and their structure is quite different from the monotone case.

\medskip
{\small
\begin{center}
\renewcommand{\arraystretch}{1.5}
\begin{tabular}{c|p{9cm}|c}
OEIS & \multicolumn{1}{c|}{Implicit rule \& denesting} & Type \\\hline
\oeis{A002516} & $a(a(n))=2n$, earliest;\;
  $a(4n)=2a(2n)$, $a(4n{+}1)=4n{+}3$, $a(4n{+}3)=8n{+}2$ & $2$-aut. \\
\oeis{A002517} & $a(a(n))=3n$, earliest;\;
  $a(3n)=3a(n)$, $a(3n{+}1)=3n{+}2$, $a(3n{+}2)=9n{+}3$ & $3$-aut. \\
\end{tabular}
\end{center}
}
\medskip

The earliest sequences take values in a rearrangement of the integers,
whereas the monotone sequences (and the almost Golomb sequences) are nondecreasing.
This difference accounts for the fact that the earliest sequences are $k$-automatic
(bounded values modulo finite-state processing) while the monotone ones are merely
$k$-regular (unbounded but finitely generated kernel).

%-------------------------------------------------------------------

\section{DFAO data for the correction sequences}\label{app:dfao}
%-------------------------------------------------------------------

This appendix provides transition tables of DFAOs
(deterministic finite automata with output) realising the correction sequences
for $r=4$ and $r=5$.
All automata read the base-$r$ expansion of $n$ from most significant to least
significant digit. The initial state is always state $0$. The tables
were exported from the Walnut certification and verified against the
recurrence systems of Theorems~\ref{thm:r4eps} and~\ref{thm:r5eps}
on a large initial range.

\emph{Note.}
The proofs in this paper do not depend on these tables.
The non-convergence results (Theorems~\ref{thm:ratio4}
and~\ref{thm:ratio5}) rely on
Lemma~\ref{lem:auto-geometric} applied to the recurrence systems
of Theorems~\ref{thm:r4eps} and~\ref{thm:r5eps},
not on the printed DFAO data.
The tables are included for reference and reproducibility only.

\subsection*{The case $r=4$}

The four correction sequences $\varepsilon_0,\varepsilon_1,\varepsilon_2,\varepsilon_3$
of Theorem~\ref{thm:r4eps} are each recognised by an explicit DFAO reading
base-$4$ expansions.

\noindent\begin{minipage}{\linewidth}
\textbf{DFAO for $\varepsilon_0$ (base $4$, 30 states, initial state $0$)}

\begin{center}
\renewcommand{\arraystretch}{1.1}
\begin{tabular}{c|c|cccc}
State & out & $d=0$ & $d=1$ & $d=2$ & $d=3$ \\\hline
0 & $0$ & 0 & 1 & 2 & 2 \\
1 & $0$ & 3 & 4 & 14 & 18 \\
2 & $0$ & 18 & 18 & 18 & 18 \\
3 & $0$ & 22 & 6 & 25 & 5 \\
4 & $0$ & 15 & 22 & 6 & 25 \\
5 & $0$ & 16 & 23 & 7 & 26 \\
6 & $0$ & 16 & 23 & 7 & 27 \\
7 & $0$ & 17 & 24 & 8 & 28 \\
8 & $0$ & 29 & 29 & 13 & 29 \\
9 & $0$ & 12 & 28 & 10 & 12 \\
10 & $0$ & 13 & 29 & 13 & 13 \\
11 & $0$ & 28 & 10 & 12 & 28 \\
12 & $0$ & 29 & 13 & 13 & 29 \\
13 & $0$ & 0 & 0 & 0 & 0 \\
14 & $1$ & 5 & 15 & 19 & 19 \\
15 & $1$ & 7 & 16 & 20 & 20 \\
16 & $1$ & 8 & 17 & 21 & 21 \\
17 & $1$ & 13 & 29 & 29 & 29 \\
18 & $1$ & 19 & 19 & 19 & 19 \\
19 & $1$ & 20 & 20 & 20 & 20 \\
20 & $1$ & 21 & 21 & 21 & 21 \\
21 & $1$ & 29 & 29 & 29 & 29 \\
22 & $1$ & 23 & 7 & 26 & 7 \\
23 & $1$ & 24 & 8 & 28 & 8 \\
24 & $1$ & 29 & 13 & 29 & 13 \\
25 & $1$ & 9 & 11 & 26 & 7 \\
26 & $1$ & 10 & 12 & 28 & 8 \\
27 & $1$ & 10 & 12 & 28 & 10 \\
28 & $1$ & 13 & 13 & 29 & 13 \\
29 & $1$ & 0 & 0 & 0 & 0 \\
\end{tabular}
\end{center}
\end{minipage}
\bigskip

\noindent\begin{minipage}{\linewidth}
\textbf{DFAO for $\varepsilon_1$ (base $4$, 30 states, initial state $0$)}

\begin{center}
\renewcommand{\arraystretch}{1.1}
\begin{tabular}{c|c|cccc}
State & out & $d=0$ & $d=1$ & $d=2$ & $d=3$ \\\hline
0 & $0$ & 0 & 6 & 1 & 1 \\
1 & $0$ & 2 & 2 & 2 & 2 \\
2 & $0$ & 3 & 3 & 3 & 3 \\
3 & $0$ & 4 & 4 & 4 & 4 \\
4 & $0$ & 5 & 5 & 5 & 5 \\
5 & $0$ & 21 & 21 & 21 & 21 \\
6 & $0$ & 9 & 22 & 15 & 2 \\
7 & $0$ & 20 & 28 & 8 & 20 \\
8 & $0$ & 21 & 29 & 21 & 21 \\
9 & $0$ & 10 & 25 & 12 & 23 \\
10 & $0$ & 11 & 26 & 13 & 24 \\
11 & $0$ & 14 & 28 & 14 & 28 \\
12 & $0$ & 19 & 27 & 13 & 24 \\
13 & $0$ & 20 & 28 & 14 & 28 \\
14 & $0$ & 21 & 29 & 21 & 29 \\
15 & $0$ & 23 & 16 & 3 & 3 \\
16 & $0$ & 24 & 17 & 4 & 4 \\
17 & $0$ & 28 & 18 & 5 & 5 \\
18 & $0$ & 29 & 21 & 21 & 21 \\
19 & $0$ & 28 & 8 & 20 & 28 \\
20 & $0$ & 29 & 21 & 21 & 29 \\
21 & $0$ & 0 & 0 & 0 & 0 \\
22 & $1$ & 16 & 10 & 25 & 12 \\
23 & $1$ & 17 & 11 & 26 & 13 \\
24 & $1$ & 18 & 14 & 28 & 14 \\
25 & $1$ & 17 & 11 & 26 & 7 \\
26 & $1$ & 18 & 14 & 28 & 8 \\
27 & $1$ & 8 & 20 & 28 & 8 \\
28 & $1$ & 21 & 21 & 29 & 21 \\
29 & $1$ & 0 & 0 & 0 & 0 \\
\end{tabular}
\end{center}
\end{minipage}
\bigskip

\noindent\begin{minipage}{\linewidth}
\textbf{DFAO for $\varepsilon_2$ (base $4$, 29 states, initial state $0$)}

\begin{center}
\renewcommand{\arraystretch}{1.1}
\begin{tabular}{c|c|cccc}
State & out & $d=0$ & $d=1$ & $d=2$ & $d=3$ \\\hline
0 & $0$ & 0 & 1 & 2 & 2 \\
1 & $0$ & 3 & 12 & 17 & 20 \\
2 & $0$ & 20 & 20 & 20 & 20 \\
3 & $0$ & 4 & 14 & 7 & 13 \\
4 & $0$ & 5 & 15 & 8 & 15 \\
5 & $0$ & 6 & 16 & 10 & 16 \\
6 & $0$ & 11 & 28 & 11 & 28 \\
7 & $0$ & 24 & 26 & 8 & 15 \\
8 & $0$ & 25 & 27 & 10 & 16 \\
9 & $0$ & 25 & 27 & 10 & 25 \\
10 & $0$ & 28 & 28 & 11 & 28 \\
11 & $0$ & 0 & 0 & 0 & 0 \\
12 & $1$ & 18 & 4 & 14 & 7 \\
13 & $1$ & 19 & 5 & 15 & 8 \\
14 & $1$ & 19 & 5 & 15 & 9 \\
15 & $1$ & 23 & 6 & 16 & 10 \\
16 & $1$ & 28 & 11 & 28 & 11 \\
17 & $1$ & 13 & 18 & 21 & 21 \\
18 & $1$ & 15 & 19 & 22 & 22 \\
19 & $1$ & 16 & 23 & 23 & 23 \\
20 & $1$ & 21 & 21 & 21 & 21 \\
21 & $1$ & 22 & 22 & 22 & 22 \\
22 & $1$ & 23 & 23 & 23 & 23 \\
23 & $1$ & 28 & 28 & 28 & 28 \\
24 & $1$ & 27 & 10 & 25 & 27 \\
25 & $1$ & 28 & 11 & 28 & 28 \\
26 & $1$ & 10 & 25 & 27 & 10 \\
27 & $1$ & 11 & 28 & 28 & 11 \\
28 & $1$ & 0 & 0 & 0 & 0 \\
\end{tabular}
\end{center}
\end{minipage}
\bigskip

\noindent\begin{minipage}{\linewidth}
\textbf{DFAO for $\varepsilon_3$ (base $4$, 30 states, initial state $0$)}

\begin{center}
\renewcommand{\arraystretch}{1.1}
\begin{tabular}{c|c|cccc}
State & out & $d=0$ & $d=1$ & $d=2$ & $d=3$ \\\hline
0 & $0$ & 0 & 5 & 1 & 1 \\
1 & $0$ & 2 & 2 & 2 & 2 \\
2 & $0$ & 3 & 3 & 3 & 3 \\
3 & $0$ & 4 & 4 & 4 & 4 \\
4 & $0$ & 9 & 9 & 9 & 9 \\
5 & $0$ & 10 & 11 & 6 & 2 \\
6 & $0$ & 12 & 7 & 3 & 3 \\
7 & $0$ & 13 & 8 & 4 & 4 \\
8 & $0$ & 16 & 9 & 9 & 9 \\
9 & $0$ & 19 & 19 & 19 & 19 \\
10 & $0$ & 20 & 14 & 22 & 12 \\
11 & $0$ & 7 & 20 & 14 & 22 \\
12 & $0$ & 8 & 21 & 15 & 23 \\
13 & $0$ & 9 & 24 & 16 & 24 \\
14 & $0$ & 8 & 21 & 15 & 25 \\
15 & $0$ & 9 & 24 & 16 & 26 \\
16 & $0$ & 19 & 29 & 19 & 29 \\
17 & $0$ & 26 & 28 & 18 & 26 \\
18 & $0$ & 29 & 29 & 19 & 29 \\
19 & $0$ & 0 & 0 & 0 & 0 \\
20 & $1$ & 21 & 15 & 23 & 13 \\
21 & $1$ & 24 & 16 & 24 & 16 \\
22 & $1$ & 27 & 17 & 23 & 13 \\
23 & $1$ & 28 & 18 & 24 & 16 \\
24 & $1$ & 29 & 19 & 29 & 19 \\
25 & $1$ & 28 & 18 & 26 & 28 \\
26 & $1$ & 29 & 19 & 29 & 29 \\
27 & $1$ & 18 & 26 & 28 & 18 \\
28 & $1$ & 19 & 29 & 29 & 19 \\
29 & $1$ & 0 & 0 & 0 & 0 \\
\end{tabular}
\end{center}
\end{minipage}
\bigskip

\subsection*{The case $r=5$}

The pair $U(n)=(\varepsilon(n),\eta(n))$ of Theorem~\ref{thm:r5eps} takes values
in $\{(0,0),(1,0),(0,1)\}$ and is recognised by an explicit DFAO reading
base-$5$ expansions (encoded as $0,1,2$ respectively).

\noindent\begin{minipage}{\linewidth}
\textbf{DFAO for $U=(\varepsilon,\eta)$ (base $5$, 23 states, initial state $0$)}

\begin{center}
\renewcommand{\arraystretch}{1.1}
\begin{tabular}{c|c|ccccc}
State & $U$ & $d=0$ & $d=1$ & $d=2$ & $d=3$ & $d=4$ \\\hline
0 & $(0,0)$ & 0 & 1 & 5 & 10 & 18 \\
1 & $(0,0)$ & 6 & 2 & 2 & 2 & 2 \\
2 & $(0,0)$ & 3 & 3 & 3 & 3 & 3 \\
3 & $(0,0)$ & 4 & 4 & 4 & 4 & 4 \\
4 & $(0,0)$ & 9 & 9 & 9 & 9 & 9 \\
5 & $(0,0)$ & 2 & 2 & 14 & 19 & 11 \\
6 & $(0,0)$ & 20 & 7 & 3 & 3 & 3 \\
7 & $(0,0)$ & 21 & 8 & 4 & 4 & 4 \\
8 & $(0,0)$ & 22 & 9 & 9 & 9 & 9 \\
9 & $(0,0)$ & 0 & 0 & 0 & 0 & 0 \\
10 & $(0,1)$ & 19 & 11 & 19 & 11 & 19 \\
11 & $(0,1)$ & 20 & 12 & 20 & 12 & 20 \\
12 & $(0,1)$ & 21 & 13 & 21 & 13 & 21 \\
13 & $(0,1)$ & 22 & 17 & 22 & 17 & 22 \\
14 & $(0,1)$ & 3 & 15 & 20 & 12 & 20 \\
15 & $(0,1)$ & 4 & 16 & 21 & 13 & 21 \\
16 & $(0,1)$ & 9 & 17 & 22 & 17 & 22 \\
17 & $(0,1)$ & 0 & 0 & 0 & 0 & 0 \\
18 & $(1,0)$ & 11 & 19 & 11 & 19 & 11 \\
19 & $(1,0)$ & 12 & 20 & 12 & 20 & 12 \\
20 & $(1,0)$ & 13 & 21 & 13 & 21 & 13 \\
21 & $(1,0)$ & 17 & 22 & 17 & 22 & 17 \\
22 & $(1,0)$ & 0 & 0 & 0 & 0 & 0 \\
\end{tabular}
\end{center}
\end{minipage}
\bigskip

\section*{Acknowledgments}

The author thanks Jean-Paul Allouche for valuable comments on the case $r=2$
and on a preliminary version of this article.

\begin{thebibliography}{9}

\bibitem{ARS05}
J.-P.\ Allouche, N.\ Rampersad and J.\ Shallit,
On integer sequences whose first iterates are linear,
\emph{Aequationes Math.}\ \textbf{69} (2005), 114--127.

\bibitem{AS92}
J.-P.\ Allouche and J.\ Shallit,
The ring of $k$-regular sequences,
\emph{Theoret.\ Comput.\ Sci.}\ \textbf{98} (1992), 163--197.
\newblock \href{https://doi.org/10.1016/0304-3975(92)90001-V}{doi:10.1016/0304-3975(92)90001-V}.

\bibitem{AS03}
J.-P.\ Allouche and J.\ Shallit,
\emph{Automatic Sequences: Theory, Applications, Generalizations},
Cambridge University Press, 2003.

\bibitem{AS12}
J.-P.\ Allouche and J.\ Shallit,
A variant of Hofstadter's sequence and finite automata,
\emph{J.\ Aust.\ Math.\ Soc.}\ \textbf{93} (2012), 1--8.
\newblock \href{https://doi.org/10.1017/S1446788712000079}{doi:10.1017/S1446788712000079}.

\bibitem{CC25}
B.\ Cloitre and J.\ Campbell,
Meta-automatic sequences,
preprint, arXiv:2602.23395 [math.CO], 2026.

\bibitem{CM26}
J.\ Claes and R.\ Miyamoto,
Golombic and Levine sequences,
preprint, arXiv:2602.10992 [math.CO], 2026.

\bibitem{Cobham}
A.\ Cobham,
Uniform tag sequences,
\emph{Math.\ Systems Theory} \textbf{6} (1972), 164--192.

\bibitem{CSV03}
B.\ Cloitre, N.~J.~A.\ Sloane, and M.~J.\ Vandermast,
Numerical analogues of Aronson's sequence,
\emph{J.\ Integer Seq.}\ \textbf{6} (2003), Article 03.2.2.

\bibitem{Golomb66}
S.~W.\ Golomb,
Problem 5407,
\emph{Amer.\ Math.\ Monthly} \textbf{73} (1966), 674.

\bibitem{Hofstadter}
D.~R.\ Hofstadter,
\emph{G\"odel, Escher, Bach: an Eternal Golden Braid},
Basic Books, 1979.

\bibitem{OEIS}
The OEIS Foundation,
\emph{The On-Line Encyclopedia of Integer Sequences},
\url{https://oeis.org}.

\bibitem{petermann95}
Y.-F.~S.\ Petermann,
On Golomb's self-describing sequence,
\emph{J.\ Number Theory} \textbf{53} (1995), 13--24.

\bibitem{Shallit2022}
J.\ Shallit,
\emph{The Logical Approach to Automatic Sequences},
Cambridge University Press, 2022.

\end{thebibliography}

\end{document}
